%% Dokumentenklasse (Koma Script) -----------------------------------------
\documentclass[%
   %draft,     % Entwurfsstadium
   final,      % fertiges Dokument
	 % --- Paper Settings ---
   paper=a4,% [Todo: add alternatives]
   paper=portrait, % landscape
   pagesize=auto, % driver
   % --- Base Font Size ---
   fontsize=11pt,%
	 % --- Koma Script Version ---
   version=last, %
 ]{scrbook} % Classes: scrartcl, scrreprt, scrbook


% Encoding der Dateien (sonst funktionieren Umlaute nicht)
% Fuer Linux -> utf8
% Fuer Windows, alte Linux Distributionen -> latin1

% Empfohlen latin1, da einige Pakete mit utf8 Zeichen nicht
% funktionieren, z.B: listings, soul.
\usepackage[latin1]{inputenc}
%\usepackage[ansinew]{inputenc}
%\usepackage[utf8]{inputenc}
%\usepackage{ucs}
%\usepackage[utf8x]{inputenc}

%%% Preambel
%%% === Textbody ==============================================================
\KOMAoptions{%
   DIV=11,% (Size of Text Body, higher values = greater textbody)
   % DIV=calc % (also areaset/classic/current/default/last) 
   % -> after setting of spacing necessary!   
   BCOR=5mm% (Bindekorrektur)
}%
%%% === Headings ==============================================================
\KOMAoptions{%
   %%%% headings
	% headings=small,  % Small Font Size, thin spacing above and below
   % headings=normal, % Medium Font Size, medium spacing above and below
   headings=big, % Big Font Size, large spacing above and below
   %
   headings=noappendixprefix, % chapter in appendix as in body text
	headings=nochapterprefix,  % no prefix at chapters
   % headings=appendixprefix,   % inverse of 'noappendixprefix'
   % headings=chapterprefix,    % inverse of 'nochapterprefix'
   % headings=openany,   % Chapters start at any side
   % headings=openleft,  % Chapters start at left side
   headings=openright, % Chapters start at right side      
   %%% Add/Dont/Auto Dot behind section numbers 
   %%% (see DUDEN as reference)
   % numbers=autoenddot
   % numbers=enddot
   numbers=noenddot
   % secnumdepth=3 % depth of sections numbering (???)
}%
\setcounter{secnumdepth}{3}
%%% === Page Layout ===========================================================
\KOMAoptions{% (most options are for package typearea)
   twoside=true, % two side layout (alternating margins, standard in books)
   % twoside=false, % single side layout 
   % twoside=semi,  % two side layout (non alternating margins!)
   %
   twocolumn=false, % (true)
   %
   headinclude=false,%
   footinclude=false,%
   mpinclude=false,%      
   %
   headlines=2.1,%
	% headheight=2em,%
   headsepline=true,%
   footsepline=false,%
   cleardoublepage=empty %plain, headings
}%
%%% === Paragraph Separation ==================================================
\KOMAoptions{%
	% parskip=relative, % change indentation according to fontsize (recommanded)
   parskip=absolute, % do not change indentation according to fontsize
   parskip=false    % indentation of 1em
   % parskip=true   % parksip of 1 line - with free space in last line of 1em
   % parskip=full-  % parksip of 1 line - no adjustment
   % parskip=full+  % parksip of 1 line - with free space in last line of 1/4
   % parskip=full*  % parksip of 1 line - with free space in last line of 1/3
   % parskip=half   % parksip of 1/2 line - with free space in last line of 1em
   % parskip=half-  % parksip of 1/2 line - no adjustment
   % parskip=half+  % parksip of 1/2 line - with free space in last line of 1/3
   % parskip=half*  % parksip of 1/2 line - with free space in last line of 1em
}%
%%% === Table of Contents =====================================================
\setcounter{tocdepth}{3} % Depth of TOC Display
\KOMAoptions{%
   %%% Setting of 'Style' and 'Content' of TOC
   % toc=left, %
   toc=indented,%
   %
   toc=bib,
   % toc=nobib,
   % toc=bibnumbered,
   %
	% toc=index,%
   toc=noindex,
   %
   % toc=listof,
   toc=nolistof
   % toc=listofnumbered,
   %   
}%  
%%% === Lists of figures, tables etc. =========================================
\KOMAoptions{%
   %%% Setting of 'Style' and 'Content' of Lists 
   %%% (figures, tables etc)
	% --- General List Style ---
   listof=left, % tabular styles
   listof=indented, % hierarchical style
   % --- chapter highlighting ---
   % listof=chapterentry, % ??? Chapter starts are marked in figure/table
   % listof=chaptergapline, % New chapter starts are marked by a gap 
      		  			   	 % of a single line
	listof=chaptergapsmall, % New chapter starts are marked by a gap 
   	    					   % of a smallsingle line
   % listof=nochaptergap, % No Gap between chapters
   %
   % listof=leveldown, % lists are moved one level down ???
   % --- Appearance of Lists in TOC
   % listof=notoc, % Lists are not part of the TOC
   listof=totoc % add Lists to TOC without number
   % listof=totocnumbered, % add Lists to TOC with number
}%  
%%% === Bibliography ==========================================================
%% Setting of 'Style' and 'Content' of Bibliography
\KOMAoptions{%
	% bibliography=oldstyle,%
   bibliography=openstyle,%
   % bibliography=nottotoc, % Bibliography is not part of the TOC
   % bibliography=totocnumbered, % add Bibliography to TOC with number
   bibliography=totoc % add Bibliography to TOC without number
}%
%%% === Index =================================================================
%% Setting of 'Style' and 'Content' of Index in TOC
\KOMAoptions{%
   index=nottotoc % index is not part of the TOC
	% index=totoc, % add index to TOC without number
}%
%%% === Titlepage =============================================================
\KOMAoptions{%
   titlepage=true %
   %titlepage=false %
}%
%%% === Miscellaneous =========================================================
\KOMAoptions{% 	
   footnotes=multiple% nomultiple
   %open=any,%
   %open=left,%
   %open=right,%
   %chapterprefix=false,%
   %appendixprefix=false,%
   %chapteratlists=10pt,% entry
}%
% ------------------------------------------------------------------------
% LaTeX - Preambel  ******************************************************
% ------------------------------------------------------------------------
% von: Matthias Pospiech
% ========================================================================

% Strukturierung dieser Praeambel:
%    1.  Pakete die vor anderen geladen werden m�ssen
%        (calc, babel, xcolor, graphicx, amsmath, pst-pdf, ragged2e, ...)
%    2.  Schriften
%    3.  Mathematik (mathtools, fixmath, onlyamsmath, braket,
%        cancel, empheq, exscale, icomma, ...)
%    4.  Tabellen (booktabs, multirow, dcolumn, tabularx, ltxtable, supertabular)
%    5.  Text
%        5.1 Auszeichnungen (ulem, soul, url)
%        5.2 Fussnoten (footmisc)
%        5.3 Verweise (varioref)
%        5.4 Listen (enumitem, paralist, declist)
%    6.  Zitieren (csquotes, jurabib, natbib)
%    7.  PDF (microtype, hyperref, backref, hypcap, pdfpages
%    8.  Graphiken (float, flafter, placeins, subfig, wrapfig,
%        floatflt, picins, psfrag, sidecap, pict2e, curve2e)
%    9.  Sonstiges (makeidx, isodate, numprint, nomencl, acronym)
%    10. Verbatim (upquote, verbatim, fancyvrb, listings, examplep)
%    11. Wissenschaft (units)
%    12. Fancy Stuff
%    13. Layout
%       13.1.  Diverse Pakete und Einstellungen (multicol, ellipsis)
%       13.2.  Zeilenabstand (setspace)
%       13.3.  Seitenlayout (typearea, geometry)
%       13.4.  Farben
%       13.5.  Aussehen der URLS
%       13.6.  Kopf und Fusszeilen (scrpage2)
%       13.7.  Fussnoten
%       13.8.  Schriften (Sections )
%       13.9.  UeberSchriften (Chapter und Sections) (titlesec, indentfirst)
%       13.10. Captions (Schrift, Aussehen)
%              (caption, subfig, capt-of, mcaption, tocloft, multitoc, minitoc)
%    14.  Auszufuehrende Befehle


% ~~~~~~~~~~~~~~~~~~~~~~~~~~~~~~~~~~~~~~~~~~~~~~~~~~~~~~~~~~~~~~~~~~~~~~~~
% Einige Pakete muessen unbedingt vor allen anderen geladen werden
% ~~~~~~~~~~~~~~~~~~~~~~~~~~~~~~~~~~~~~~~~~~~~~~~~~~~~~~~~~~~~~~~~~~~~~~~~

%%% Packages for LaTeX - programming
%
% Define commands that don't eat spaces.
\usepackage{xspace}
% IfThenElse
\usepackage{ifthen}
%%% Doc: ftp://tug.ctan.org/pub/tex-archive/macros/latex/contrib/oberdiek/ifpdf.sty
% command for testing for pdf-creation
\usepackage{ifpdf} %\ifpdf \else \fi

%%% Internal Commands: ----------------------------------------------
\makeatletter
%
\providecommand{\IfPackageLoaded}[2]{\@ifpackageloaded{#1}{#2}{}}
\providecommand{\IfPackageNotLoaded}[2]{\@ifpackageloaded{#1}{}{#2}}
\providecommand{\IfElsePackageLoaded}[3]{\@ifpackageloaded{#1}{#2}{#3}}
%
\newboolean{chapteravailable}%
\setboolean{chapteravailable}{false}%

\ifcsname chapter\endcsname
  \setboolean{chapteravailable}{true}%
\else
  \setboolean{chapteravailable}{false}%
\fi


\providecommand{\IfChapterDefined}[1]{\ifthenelse{\boolean{chapteravailable}}{#1}{}}%
\providecommand{\IfElseChapterDefined}[2]{\ifthenelse{\boolean{chapteravailable}}{#1}{#2}}%

\providecommand{\IfDefined}[2]{%
\ifcsname #1\endcsname
   #2 %
\else
     % do nothing
\fi
}

\providecommand{\IfElseDefined}[3]{%
\ifcsname #1\endcsname
   #2 %
\else
   #3 %
\fi
}

\providecommand{\IfElseUnDefined}[3]{%
\ifcsname #1\endcsname
   #3 %
\else
   #2 %
\fi
}


%
% Check for 'draft' mode - commands.
\newcommand{\IfNotDraft}[1]{\ifx\@draft\@undefined #1 \fi}
\newcommand{\IfNotDraftElse}[2]{\ifx\@draft\@undefined #1 \else #2 \fi}
\newcommand{\IfDraft}[1]{\ifx\@draft\@undefined \else #1 \fi}
%

% Definde frontmatter, mainmatter and backmatter if not defined
\@ifundefined{frontmatter}{%
   \newcommand{\frontmatter}{%
      %In Roemischen Buchstaben nummerieren (i, ii, iii)
      \pagenumbering{roman}
   }
}{}
\@ifundefined{mainmatter}{%
   % scrpage2 benoetigt den folgenden switch
   % wenn \mainmatter definiert ist.
   \newif\if@mainmatter\@mainmattertrue
   \newcommand{\mainmatter}{%
      % -- Seitennummerierung auf Arabische Zahlen zuruecksetzen (1,2,3)
      \pagenumbering{arabic}%
      \setcounter{page}{1}%
   }
}{}
\@ifundefined{backmatter}{%
   \newcommand{\backmatter}{
      %In Roemischen Buchstaben nummerieren (i, ii, iii)
      \pagenumbering{roman}
   }
}{}

% Pakete speichern die spaeter geladen werden sollen
\newcommand{\LoadPackagesNow}{}
\newcommand{\LoadPackageLater}[1]{%
   \g@addto@macro{\LoadPackagesNow}{%
      \usepackage{#1}%
   }%
}



\makeatother
%%% ----------------------------------------------------------------
%
%%% Doc: www.cs.brown.edu/system/software/latex/doc/calc.pdf
% Calculation with LaTeX
\usepackage{calc}

%%% Doc: ftp://tug.ctan.org/pub/tex-archive/macros/latex/required/babel/babel.pdf
% Languagesetting
\usepackage[
%	german,
%	ngerman,
	english,
%	french,
]{babel}

%%% Doc: ftp://tug.ctan.org/pub/tex-archive/macros/latex/contrib/xcolor/xcolor.pdf
% Farben
% Incompatible: Do not load when using pstricks !
\usepackage[
	table % Load for using rowcolors command in tables
]{xcolor}


%%% Doc: ftp://tug.ctan.org/pub/tex-archive/macros/latex/required/graphics/grfguide.pdf
% Bilder
\usepackage[%
	%final,
	%draft % do not include images (faster)
]{graphicx}

%%% Doc: ftp://tug.ctan.org/pub/tex-archive/macros/latex/contrib/oberdiek/epstopdf.pdf
%% If an eps image is detected, epstopdf is automatically called to convert it to pdf format.
%% Requires: graphicx loaded
\usepackage{epstopdf}


%%% Doc: ftp://tug.ctan.org/pub/tex-archive/macros/latex/required/amslatex/math/amsldoc.pdf
% Amsmath - Mathematik Basispaket
%
% fuer pst-pdf displaymath Modus vor pst-pdf benoetigt.
\usepackage[
   centertags, % (default) center tags vertically
   %tbtags,    % 'Top-or-bottom tags': For a split equation, place equation numbers level
               % with the last (resp. first) line, if numbers are on the right (resp. left).
   sumlimits,  %(default) Place the subscripts and superscripts of summation
               % symbols above and below
   %nosumlimits, % Always place the subscripts and superscripts of summation-type
               % symbols to the side, even in displayed equations.
   intlimits,  % Like sumlimits, but for integral symbols.
   %nointlimits, % (default) Opposite of intlimits.
   namelimits, % (default) Like sumlimits, but for certain 'operator names' such as
               % det, inf, lim, max, min, that traditionally have subscripts placed underneath
               % when they occur in a displayed equation.
   %nonamelimits, % Opposite of namelimits.
   %leqno,     % Place equation numbers on the left.
   %reqno,     % Place equation numbers on the right.
   fleqn,     % Position equations at a fixed indent from the left margin
   			  % rather than centered in the text column.
]{amsmath} %
% eqnarray nicht zusammen mit amsmath benutzen, siehe l2tabu.pdf f�r
% Hintergruende.

%%% Doc: http://www.ctan.org/tex-archive/macros/latex/contrib/pst-pdf/pst-pdf-DE.pdf
% Used to automatically integrate eps graphics in an pdf document using pdflatex.
% Requires ps4pdf macro !!!
% Download macro from http://www.ctan.org/tex-archive/macros/latex/contrib/pst-pdf/scripts/
%
%\usepackage[%
%   %active,       % Aktiviert den Extraktionsmodus (DVI-Ausgabe). Die explizite Angabe ist
%                  % normalerweise unn�tig (Standard im LATEX-Modus).
%   %inactive,     % Das Paket wird deaktiviert, Zu�tzlich werden die Pakete pstricks und
%                  % graphicx geladen
%   nopstricks,    % Das Paket pstricks wird nicht geladen.
%   %draft,        % Im pdfLATEX-Modus werden aus der Containerdatei eingef�gte Grafiken nur
%                  % als Rahmen dargestellt.
%   %final,        % Im pdfLATEX-Modus werden aus der Containerdatei eingef�gte Grafiken
%                  % vollst�ndig dargestellt (Standard).
%   %tightpage,    % Die Abmessung Grafiken in der Containerdatei entsprechen denen der
%                  % zugeh�rigen TEX-Boxen (Standard).
%   %notightpage,  % die Grafiken in der Containerdatei nehmen
%                  % mindestens die Gr��e des gesamten Blattes einnehmen.
%   %displaymath,  % Es werden zus�tzlich die mathematischen Umgebungen displaymath,
%                  % eqnarray und $$ extrahiert und im pdf-Modus als Grafik eingef�gt.
%]{pst-pdf}
%
% Notwendiger Bugfix f�r natbib Paket bei Benutzung von pst-pdf (Version <= v1.1o)
\IfPackageLoaded{pst-pdf}{
   \providecommand\makeindex{}
   \providecommand\makeglossary{}
}{}


%% Doc: ftp://tug.ctan.org/pub/tex-archive/graphics/pstricks/README
% load before graphicx
% \usepackage{pstricks}
% \usepackage{pst-plot, pst-node, pst-coil, pst-eps}

% This package implements a workaround for the LaTeX bug that marginpars
% sometimes appear on the wrong margin.
% \usepackage{mparhack}
% in some case this causes an error in the index together with package pdfpages
% the reason is unkown. Therefore I recommend to use the margins of marginnote

%% Doc: ftp://tug.ctan.org/pub/tex-archive/macros/latex/contrib/marginnote/marginnote.pdf
% Summary description: marginnote allows margin note, where \marginpar fails
\usepackage{marginnote}


%% Doc: (inside relsize.sty )
%% ftp://tug.ctan.org/pub/tex-archive/macros/latex/contrib/misc/relsize.sty
%  Set the font size relative to the current font size
\usepackage{relsize}

%% Doc: ftp://tug.ctan.org/pub/tex-archive/macros/latex/contrib/ms/ragged2e.pdf
% Besserer Flatternsatz (Linksbuendig, statt Blocksatz)
\usepackage{ragged2e}

% ~~~~~~~~~~~~~~~~~~~~~~~~~~~~~~~~~~~~~~~~~~~~~~~~~~~~~~~~~~~~~~~~~~~~~~~~
% Fonts Fonts Fonts
% ~~~~~~~~~~~~~~~~~~~~~~~~~~~~~~~~~~~~~~~~~~~~~~~~~~~~~~~~~~~~~~~~~~~~~~~~

\usepackage[T1]{fontenc} % T1 Schrift Encoding
\usepackage{textcomp}	 % Zusatzliche Symbole (Text Companion font extension)

%%% Schriften werden in Fonts.tex geladen
% ~~~~~~~~~~~~~~~~~~~~~~~~~~~~~~~~~~~~~~~~~~~~~~~~~~~~~~~~~~~~~~~~~~~~~~~~
% Fonts Fonts Fonts
% ~~~~~~~~~~~~~~~~~~~~~~~~~~~~~~~~~~~~~~~~~~~~~~~~~~~~~~~~~~~~~~~~~~~~~~~~

% Alle Schriften die hier angegeben sind sehen im PDF richtig aus.
% Die LaTeX Standardschrift ist die Latin Modern (lmodern Paket).
% If Latin Modern is not available for your distribution you must install the
% package cm-super instead. Otherwise your fonts will look horrible in the PDF

% DO NOT LOAD ae Package for the font !

%% ==== Zusammengesetzte Schriften  (Sans + Serif) =======================

%% - Latin Modern
\usepackage{lmodern}
%% -------------------
%
%% - Times, Helvetica, Courier (Word Standard...)
%\usepackage{mathptmx}
%\usepackage[scaled=.90]{helvet}
%\usepackage{courier}
%% -------------------
%%
%% - Palantino , Helvetica, Courier
%\usepackage{mathpazo}
%\usepackage[scaled=.95]{helvet}
%\usepackage{courier}
%% -------------------
%
%% - Bera Schriften
%\usepackage{bera}
%% -------------------
%
%% - Charter, Bera Sans
%\usepackage{charter}\linespread{1.05}
%\renewcommand{\sfdefault}{fvs}

%% ===== Serifen =========================================================

%\usepackage{mathpazo}                 %% --- Palantino
%\usepackage{charter}\linespread{1.05} %% --- Charter
%\usepackage{bookman}                  %% --- Bookman (laedt Avant Garde !!)
%\usepackage{newcent}                  %% --- New Century Schoolbook (laedt Avant Garde !!)

%\usepackage[%                         %% --- Fourier
%   upright,     % Math fonts are upright
%   expert,      % Only for EXPERT Fonts!
%   oldstyle,    % Only for EXPERT Fonts!
%   fulloldstyle % Only for EXPERT Fonts!
%]{fourier} %



%% ===== Sans Serif ======================================================

%\usepackage[scaled=.95]{helvet}        %% --- Helvetica
%\usepackage{cmbright}                  %% --- CM-Bright (eigntlich eine Familie)
%\usepackage{tpslifonts}                %% --- tpslifonts % Font for Slides
%\usepackage{avantgar}                  %% --- Avantgarde

%%%% =========== Italics ================

%\usepackage{chancery}                  %% --- Zapf Chancery

%%%% =========== Typewriter =============

%\usepackage{courier}                   %% --- Courier
%\renewcommand{\ttdefault}{cmtl}        %% --- CmBright Typewriter Font
%\usepackage[%                          %% --- Luxi Mono (Typewriter)
%   scaled=0.9
%]{luximono}



%%%% =========== Mathe ================

%% Recommanded to use with fonts: Aldus, Garamond, Melior, Sabon
%\usepackage[                           %% --- EulerVM (MATH)
%   small,       %for smaller Fonts
%  euler-digits % digits in euler fonts style
%]{eulervm}

% \usepackage[
% %   utopia,
% %   garamond,
%    charter
% ]{mathdesign}

%%%% (((( !!! kommerzielle Schriften !!! )))))))))))))))))))))))))))))))))))))))))))))))))))

%% ===== Serifen (kommerzielle Schriften ) ================================

%% --- Adobe Aldus
%\renewcommand{\rmdefault}{pasx}
%\renewcommand{\rmdefault}{pasj} %%oldstyle digits
% math recommended: \usepackage[small]{eulervm}

%% --- Adobe Garamond
%\usepackage[%
%   osf,        % oldstyle digits
%   scaled=1.05 %appropriate in many cases
%]{xagaramon}
% math recommended: \usepackage{eulervm}

%% --- Adobe Stempel Garamond
%\renewcommand{\rmdefault}{pegx}
%\renewcommand{\rmdefault}{pegj} %%oldstyle digits

%% --- Adobe Melior
%\renewcommand{\rmdefault}{pml}
% math recommended: %\usepackage{eulervm}

%% --- Adobe Minion
%\renewcommand{\rmdefault}{pmnx}
%\renewcommand{\rmdefault}{pmnj} %oldstyle digits
% math recommended: \usepackage[small]{eulervm} or \usepackage{mathpmnt} % commercial

%% --- Adobe Sabon
%\renewcommand{\rmdefault}{psbx}
%\renewcommand{\rmdefault}{psbj} %oldstyle digits
% math recommended: \usepackage{eulervm}

%% --- Adobe Times
% math recommended: \usepackage{mathptmx} % load first !
%\renewcommand{\rmdefault}{ptmx}
%\renewcommand{\rmdefault}{ptmj} %oldstyle digits

%% --- Linotype ITC Charter
%\renewcommand{\rmdefault}{lch}

%% --- Linotype Meridien
%\renewcommand{\rmdefault}{lmd}

%%% ===== Sans Serif (kommerzielle Schriften) ============================

%% --- Adobe Frutiger
%\usepackage[
%   scaled=0.90
%]{frutiger}

%% --- Adobe Futura (=Linotype FuturaLT) : Sans Serif
%\usepackage[
%   scaled=0.94  % appropriate in many cases
%]{futura}

%% --- Adobe Gill Sans : Sans Serif
%\usepackage{gillsans}

%% -- Adobe Myriad  : Sans Serif
%\renewcommand{\sfdefault}{pmy}
%\renewcommand{\sfdefault}{pmyc} %% condensed Font

%% --- Syntax : sans serif font
%\usepackage[
%   scaled
%]{asyntax}

%% --- Adobe Optima : Semi Sans Serif
%\usepackage[
%   medium %darker medium weight fonts
%]{optima}

%% --- Linotype ITC Officina Sans
%\renewcommand{\sfdefault}{lo9}





% ~~~~~~~~~~~~~~~~~~~~~~~~~~~~~~~~~~~~~~~~~~~~~~~~~~~~~~~~~~~~~~~~~~~~~~~~
% Math Packages
% ~~~~~~~~~~~~~~~~~~~~~~~~~~~~~~~~~~~~~~~~~~~~~~~~~~~~~~~~~~~~~~~~~~~~~~~~

% *** Mathematik **************************************
%
% amsmath schon vorher geladen da es vor pst-pdf geladen werden muss


%%% Doc: ftp://tug.ctan.org/pub/tex-archive/macros/latex/contrib/mh/doc/mathtools.pdf
% Erweitert amsmath und behebt einige Bugs
\usepackage[fixamsmath,disallowspaces]{mathtools}

%%% Doc: http://www.ctan.org/info?id=fixmath
% LaTeX's default style of typesetting mathematics does not comply
% with the International Standards ISO31-0:1992 to ISO31-13:1992
% which indicate that uppercase Greek letters always be typset
% upright, as opposed to italic (even though they usually
% represent variables) and allow for typsetting of variables in a
% boldface italic style (even though the required fonts are
% available). This package ensures that uppercase Greek be typeset
% in italic style, that upright $\Delta$ and $\Omega$ symbols are
% available through the commands \upDelta and \upOmega; and
% provides a new math alphabet \mathbold for boldface
% italic letters, including Greek.
\usepackage{fixmath}

%%% Doc: ftp://tug.ctan.org/pub/tex-archive/macros/latex/contrib/onlyamsmath/onlyamsmath.dvi
% Warnt bei Benutzung von Befehlen die mit amsmath inkompatibel sind.
\usepackage[
	all,
	warning
]{onlyamsmath}


%------------------------------------------------------

% -- Vektor fett darstellen -----------------
% \let\oldvec\vec
% \def\vec#1{{\boldsymbol{#1}}} %Fetter Vektor
% \newcommand{\ve}{\vec} %
% -------------------------------------------


%%% Doc: ftp://tug.ctan.org/pub/tex-archive/macros/latex/contrib/misc/braket.sty
\usepackage{braket}  % Quantenmechanik Bracket Schreibweise

%%% Doc: ftp://tug.ctan.org/pub/tex-archive/macros/latex/contrib/misc/cancel.sty
\usepackage{cancel}  % Durchstreichen

%%% Doc: ftp://tug.ctan.org/pub/tex-archive/macros/latex/contrib/mh/doc/empheq.pdf
\usepackage{empheq}  % Hervorheben

%%% Doc: ftp://tug.ctan.org/pub/tex-archive/info/math/voss/mathmode/Mathmode.pdf
%\usepackage{exscale} % Skaliert Mathe-Modus Ausgaben in allen Umgebungen richtig.

%%% Doc: ftp://tug.ctan.org/pub/tex-archive/macros/latex/contrib/was/icomma.dtx
% Erlaubt die Benutzung von Kommas im Mathematikmodus
\usepackage{icomma}


%%% Doc: http://www.ctex.org/documents/packages/special/units.pdf
% \usepackage[nice]{nicefrac}

%%% Tauschen von Epsilon und andere:
% \let\ORGvarrho=\varrho
% \let\varrho=\rho
% \let\rho=\ORGvarrho
%
\let\ORGvarepsilon=\varepsilon
\let\varepsilon=\epsilon
\let\epsilon=\ORGvarepsilon
%
% \let\ORGvartheta=\vartheta
% \let\vartheta=\theta
% \let\theta=\ORGvartheta
%
% \let\ORGvarphi=\varphi
% \let\varphi=\phi
% \let\phi=\ORGvarphi

% ~~~~~~~~~~~~~~~~~~~~~~~~~~~~~~~~~~~~~~~~~~~~~~~~~~~~~~~~~~~~~~~~~~~~~~~~
% Symbole
% ~~~~~~~~~~~~~~~~~~~~~~~~~~~~~~~~~~~~~~~~~~~~~~~~~~~~~~~~~~~~~~~~~~~~~~~~
%
%%% General Doc: http://www.ctan.org/tex-archive/info/symbols/comprehensive/symbols-a4.pdf
%
%% Symbole f�r Mathematiksatz
%\usepackage{mathrsfs} %% Schreibschriftbuchstaben f�r den Mathematiksatz (nur Gro�buchstaben)
%\usepackage{dsfont}   %% Double Stroke Fonts
%\usepackage[mathcal]{euscript} %% adds euler mathcal font
%\usepackage{amssymb}
\usepackage[Symbolsmallscale]{upgreek} % upright symbols from euler package [Euler] or Adobe Symbols [Symbol]
\usepackage[upmu]{gensymb}             % Option upmu

%% Allgemeine Symbole
%\usepackage{wasysym}  %% Doc: http://www.ctan.org/tex-archive/macros/latex/contrib/wasysym/wasysym.pdf
%\usepackage{marvosym} %% Symbole aus der marvosym Schrift
%\usepackage{pifont}   %% ZapfDingbats


% ~~~~~~~~~~~~~~~~~~~~~~~~~~~~~~~~~~~~~~~~~~~~~~~~~~~~~~~~~~~~~~~~~~~~~~~~
% Tables (Tabular)
% ~~~~~~~~~~~~~~~~~~~~~~~~~~~~~~~~~~~~~~~~~~~~~~~~~~~~~~~~~~~~~~~~~~~~~~~~

% Basispaket fuer alle Tabellenfunktionen
% -> wird automatisch durch andere Pakete geladen
% \usepackage{array}
%
% bessere Abstaende innerhalb der Tabelle (Layout))
% -------------------------------------------------
%%% Doc: ftp://tug.ctan.org/pub/tex-archive/macros/latex/contrib/booktabs/booktabs.pdf
\usepackage{booktabs}
%
% Farbige Tabellen
% ----------------
% Das Paket colortbl wird inzwischen automatisch durch xcolor geladen
%
% Erweiterte Funktionen innerhalb von Tabellen
% --------------------------------------------
%%% Doc: ftp://tug.ctan.org/pub/tex-archive/macros/latex/contrib/multirow/multirow.sty
\usepackage{multirow} % Mehrfachspalten
%
%%% Doc: Documentation inside dtx Package
\usepackage{dcolumn}  % Ausrichtung an Komma oder Punkt

%%% Neue Tabellen-Umgebungen:
% ---------------------------
% Spalten automatischer Breite:
%%% Doc: Documentation inside dtx Package
% \usepackage{tabularx}
% -> nach hyperref Laden
% -> wird von ltxtable geladen
% \LoadPackageLater{tabularx}


% Tabellen ueber mehere Seiten
% ----------------------------
%%% Doc: ftp://tug.ctan.org/pub/tex-archive/macros/latex/contrib/carlisle/ltxtable.pdf
% \usepackage{ltxtable} % Longtable + tabularx
                        % (multi-page tables) + (auto-sized columns in a fixed width table)
% -> nach hyperref laden
\LoadPackageLater{ltxtable}


%%% Doc: ftp://tug.ctan.org/pub/tex-archive/macros/latex/contrib/supertabular/supertabular.pdf
%\usepackage{supertabular}


% ~~~~~~~~~~~~~~~~~~~~~~~~~~~~~~~~~~~~~~~~~~~~~~~~~~~~~~~~~~~~~~~~~~~~~~~~
% text related packages
% ~~~~~~~~~~~~~~~~~~~~~~~~~~~~~~~~~~~~~~~~~~~~~~~~~~~~~~~~~~~~~~~~~~~~~~~~

%%% Textverzierungen/Auszeichnungen ======================================
%
%%% Doc: ftp://tug.ctan.org/pub/tex-archive/macros/latex/contrib/misc/ulem.sty
\usepackage[normalem]{ulem}      % Zum Unterstreichen
%%% Doc: ftp://tug.ctan.org/pub/tex-archive/macros/latex/contrib/soul/soul.pdf
\usepackage{soul}		            % Unterstreichen, Sperren
%%% Doc: ftp://tug.ctan.org/pub/tex-archive/macros/latex/contrib/misc/url.sty
\usepackage{url} % Setzen von URLs. In Verbindung mit hyperref sind diese auch aktive Links.

%%% Fussnoten/Endnoten ===================================================
%
%%% Doc: ftp://tug.ctan.org/pub/tex-archive/macros/latex/contrib/footmisc/footmisc.pdf
%
\usepackage[
   bottom,      % Footnotes appear always on bottom. This is necessary
                % especially when floats are used
   stable,      % Make footnotes stable in section titles
   perpage,     % Reset on each page
   %para,       % Place footnotes side by side of in one paragraph.
   %side,       % Place footnotes in the margin
   ragged,      % Use RaggedRight
   %norule,     % suppress rule above footnotes
   multiple,    % rearrange multiple footnotes intelligent in the text.
   %symbol,     % use symbols instead of numbers
]{footmisc}

\renewcommand*{\multfootsep}{,\nobreakspace}

\deffootnote%
   [1em]% width of marker
   {1.5em}% indentation (general)
   {1em}% indentation (par)
   {\textsubscript{\thefootnotemark}}%


%% Einruecken der Fussnote einstellen
%\setlength\footnotemargin{10pt}

%--- footnote counter documentweit durchlaufend ------------------------------
%\usepackage{chngcntr}
%\counterwithout{footnote}{chapter}
%-----------------------------------------------------------------------------

%%% Doc: ftp://tug.ctan.org/pub/tex-archive/macros/latex/contrib/misc/endnotes.sty
%\usepackage{endnotes}
% From the Documentation:
% To turn all the footnotes in your documents into endnotes, say
%
%     \let\footnote=\endnote
%
%  in your preamble, and then add something like
%
%     \newpage
%     \begingroup
%     \parindent 0pt
%     \parskip 2ex
%     \def\enotesize{\normalsize}
%     \theendnotes
%     \endgroup
%
% as the last thing in your document.  (But \theendnotes all
% by itself will work.)

%%% Verweise =============================================================
%
%%% Doc: Documentation inside dtx File
\usepackage[english]{varioref} % Intelligente Querverweise

%%% Listen ===============================================================
%
%
%%% Doc: ftp://tug.ctan.org/pub/tex-archive/macros/latex/contrib/paralist/paralist.pdf
% \usepackage{paralist}
%
%%% Doc: ftp://tug.ctan.org/pub/tex-archive/macros/latex/contrib/enumitem/enumitem.pdf
% Better than 'paralist' and 'enumerate' because it uses a keyvalue interface !
% Do not load together with enumerate.
\IfPackageNotLoaded{enumerate}{
	\usepackage{enumitem}
}
%
%%% Doc: ftp://tug.ctan.org/pub/tex-archive/macros/latex/contrib/ncctools/doc/desclist.pdf
% Improved description environment
%\usepackage{declist}


% ~~~~~~~~~~~~~~~~~~~~~~~~~~~~~~~~~~~~~~~~~~~~~~~~~~~~~~~~~~~~~~~~~~~~~~~~
% Pakete zum Zitieren
% ~~~~~~~~~~~~~~~~~~~~~~~~~~~~~~~~~~~~~~~~~~~~~~~~~~~~~~~~~~~~~~~~~~~~~~~~

% Quotes =================================================================
%% Doc: ftp://tug.ctan.org/pub/tex-archive/macros/latex/contrib/csquotes/csquotes.pdf
% Advanced features for clever quotations
\usepackage[%
   babel,            % the style of all quotation marks will be adapted
                     % to the document language as chosen by 'babel'
   german=quotes,		% Styles of quotes in each language
   english=british,
   french=guillemets
]{csquotes}

% All facilities which take a 'cite' argument will not insert
% it directly. They pass it to an auxiliary command called \mkcitation
% which  may be redefined to format the citation.
\renewcommand*{\mkcitation}[1]{{\,}#1}
\renewcommand*{\mkccitation}[1]{ #1}

\SetBlockThreshold{2} % Anzahl von Zeilen

\newenvironment{myquote}%
	{\begin{quote}\small}%
	{\end{quote}}%
\SetBlockEnvironment{myquote}
%\SetCiteCommand{} % Changes citation command


% Zitate =================================================================
%
% Reference by number
% \makeatletter
% \renewcommand\@biblabel[1]{#1.}
% \makeatother

% -------- Reference by author -------------------------------------------

% Doc: http://www.berger-on.net/jurabib/
% \usepackage{jurabib}
% \jurabibsetup{
%    authorformat={
%       %abbrv,%  First names will be abbreviated
%       %allreversed,% Names will be printed reversed ('first surname' in text and bibliography)
%       %citationreversed,% Names will be printed reversed ('first surname' only in text)
%       and,% Author separation with ',' and ', and' instead of the default slashes
%       %dynamic,% Font of author depends on existence of coauthor
%       %firstnotreversed, %All author names (except the first) will be printed reversed (in text only)
%       %indexed,% All author names are indexed separately (makeidx or index have to be loaded appropriately)
%       %italic,% Author will appear in italics
%       %reducedifibidem,%  Author names are reduced to the surname for subsequent
%       %citations,% (for the ibidem=name options)
%       smallcaps,% Author will appear in small caps
%       year,% Emulates author-year citations
%    },
%    bibformat={
%       %compress,% Reduces the vertical space between the items in the bibliography
%       %ibidem,%  Replaces repeated authors in the bibliography
%       %ibidemalt,% Special format for German law students
%       %nohang,% No hanging indent for bibliography
%       numbered,% Numbered items in bibliography
%       %raggedright,% Flushleft for bibliography
%       %tabular,% Tabular-like bibliography format
%    },
%    titleformat={
%       all,% Prints out all titles, doesn't care about multiple works
%       %colonsep,% Separation between author and title with colon
%       comma,% Separation between author and title with comma
%       italic,% Title will appear in italics
%    },
%    %biblikecite=true,% Formatting of bibliography follows formatting of citations (as far as possible)
%    %coauthorformat={
%    %  normal,% No special format for coauthors
%    %  %italic,% Coauthor will appear in italics
%    %},
%    %colastsep=divis,% Coauthor after author, separation with divis, Standard is a slash
%    %cofirstsep={
%    %% in,% Author after coauthor, separation with 'in'
%    %  comma,% Author after coauthor, separation with comma
%    %},
%    commabeforerest,  % Nach allen Angaben und vor den zusaetzlichen
%                      % (z.B. Seitenanzahl) wird ein Komma gesetzt
%    citefull={
%       %all,% All citations are full citations
%       first,% First citation is printed full
%       %chapter,% citefull=first, resetted each chapter (book and report classes)
%       %section,% citefull=first, resetted each section (article classes)
%    },
%    %chicago=true, %chicago-like format of citation and bibliography
%    %oxford=true, %Emulates oxford-like format of citation and bibliography
%    %crossref={
%    %  dynamic,% Long crossref's if they are used first time, shorter for all further citations
%    %  long,% Always long crossref's
%    %  short,% Always short crossref's (short as possible, longer if citations are ambiguous)
%    %},
%    %edby=true,% Switches from '(ed.)' to 'edited by' for incollections
%    %endnote=true,% The note field is printed at the end of the entry, after the closing period
%    %footnotes=marginal,% Another footnote format
%    %human=true, % Common humanities option, make authorformat=and the default
%    %howcited={
%    %  all,% The howcited remark is printed for all entries
%    %  compare,% The howcited remark is printed for works, where title and shorttitle differ
%    %  normal,% The howcited remark is printed for entries containing a non-empty howcited field
%    %  multiple,%  The howcited remark is printed if more than one work of the author is cited
%    %},
%    ibidem={
%       %name,% Ibidem with authors name
%       %name&title,% Ibidem with authors name and title
%       %nostrict,% Ibidem is allowed for every footnote
%       strict,% Ibidem is not allowed for first footnote on each page
%       %strictdoublepage,% Ibidem is not allowed for first footnote on left (even) pages
%    },
%    idem={
%       %nostrict,% Idem is allowed for every footnote
%       strict,% Idem is not allowed for first footnote on each page
%       %strictdoublepage,%  Idem is not allowed for first footnote on left (even) pages
%    },
%    %lookat=true,% Enables crossref to full (first) footnote citation
%    %natoptargorder=true,% Reversed optional arguments
%    %opcit={
%    %  true,% Enables op. cit. for already cited, but not subsequent works
%    %  chapter,% Resets opcit each chapter (book and report classes)
%    %  section,% Resets opcit each section (article classes)
%    %},
%    pages={
%       %always,%  Page(ranges)s given via the pages-field are always printed in the citation
%       format,%  Pages given via the optional argument and page(ranges)s given via the pages-field are formatted automatically
%       %test,% Page(ranges)s given via the pages-field are printed in the citation, if no pages are given by the optional argument
%    },
%    %see=true,% The second optional argument can be used to add sequences like 'see' before the citation
%    %superscriptedition={
%    %  all,% Superscripted edition number for all citations
%    %  bib,% Superscripted edition number for the bibliography
%    %  commented,% Superscripted edition number only for type @COMMENTED
%    %  switch,%  Superscripted edition number for works with field ssedition=1
%    %},
%    super, % alle \cite werden zu footcite
% }
%
% \bibliographystyle{jurabib} %
% %\bibliographystyle{jureco}
% %\bibliographystyle{jurunsrt}
% %\bibliographystyle{jox}
%
% \renewcommand{\biblnfont}{\bfseries\scshape\RaggedRight} % Autoren
% \renewcommand{\bibfnfont}{} % Autoren Vornamen
% \renewcommand{\bibelnfont}{\normalfont} % Herausgeber
% \renewcommand{\bibefnfont}{\normalfont} % Herausgeber Vornamen
% % Anpassung der Titel von Buechern
% \renewcommand{\bibtfont}{\normalfont\textit} % Titel
% % % Modifizierung des Zeitschriftentitels bei Artikeln.
% \renewcommand{\bibjtfont}{\normalfont}
% % % Titel eines Artikels, eines Beitrages in einem Sammelwerk oder aehnliches zu formatieren.
% \renewcommand{\bibapifont}{\normalfont\textit} % Periodical-Titel
% % % Aussehen des series Feldes bestimmen
% \renewcommand{\bibsnfont}{\textbf}
% %
% \renewcommand{\biburlprefix}{Webseite:{ }}
% \renewcommand{\biburlsuffix}{}
%
% % Ausgabe des Jahres
% \renewcommand{\jbcitationyearformat}[1]{(#1)}
%
% \renewcommand{\bibleftcolumnadjust}{\RaggedRight}
% \renewcommand{\bibrightcolumnadjust}{\RaggedRight}
% %
% -------- Reference by number (author) ----------------------------------

% %%% Doc: ftp://tug.ctan.org/pub/tex-archive/macros/latex/contrib/natbib/natbib.pdf
%\usepackage[%
%	%round,	%(default) for round parentheses;
%	square,	% for square brackets;
%	%curly,	% for curly braces;
%	%angle,	% for angle brackets;
%	%colon,	% (default) to separate multiple citations with colons;
%	comma,	% to use commas as separaters;
%	%authoryear,% (default) for author-year citations;
%	numbers,	% for numerical citations;
%	%super,	% for superscripted numerical citations, as in Nature;
%	sort,		% orders multiple citations into the sequence in which they appear in the list of references;
%	sort&compress,    % as sort but in addition multiple numerical citations
%                   % are compressed if possible (as 3-6, 15);
%	%longnamesfirst,  % makes the first citation of any reference the equivalent of
%                   % the starred variant (full author list) and subsequent citations
%                   %normal (abbreviated list);
%	%sectionbib,      % redefines \thebibliography to issue \section* instead of \chapter*;
%                   % valid only for classes with a \chapter command;
%                   % to be used with the chapterbib package;
%	%nonamebreak,     % keeps all the authors names in a citation on one line;
%                   %causes overfull hboxes but helps with some hyperref problems.
%]{natbib}

%\RequirePackage[fixlanguage]{babelbib}
%\bibliographystyle{babalpha-fl}%

%%% Bibliography styles with natbib support
%\bibliographystyle{plainnat} % Numeric Labels, alphabatical order
%\bibliographystyle{abbrvnat} % same as plain, but shorter names
%\bibliographystyle{unsrtnat} % same as plain, but appeariance in order of citation
%\bibliographystyle{alpha}    % labels are formed by author and year

%%% Bibliography styles according to DIN
%%% get from: http://www.ctan.org/tex-archive/biblio/bibtex/contrib/german/din1505/
%\bibliographystyle{alphadin}
%\bibliographystyle{abbrvdin}
%\bibliographystyle{plaindin}
%\bibliographystyle{unsrtdin}
%\bibliographystyle{bib/bst/alphadin-mod} % Modifiziert: Kleinere Abstaende vor ";" und kein "+" bei etal.

%%% Bibliography styles created with custombib
%%% Doc: ftp://tug.ctan.org/pub/tex-archive/macros/latex/contrib/custom-bib/makebst.pdf
\bibliographystyle{bib/bst/AlphaDINFirstName}

% other BibTeX styles: http://www.cs.stir.ac.uk/~kjt/software/latex/showbst.html


% ~~~~~~~~~~~~~~~~~~~~~~~~~~~~~~~~~~~~~~~~~~~~~~~~~~~~~~~~~~~~~~~~~~~~~~~~
% figures and placement
% ~~~~~~~~~~~~~~~~~~~~~~~~~~~~~~~~~~~~~~~~~~~~~~~~~~~~~~~~~~~~~~~~~~~~~~~~

%% Bilder und Graphiken ==================================================

%%% Doc: only dtx Package
\usepackage{float}             % Stellt die Option [H] fuer Floats zur Verfgung

%%% Doc: No Documentation
\usepackage{flafter}          % Floats immer erst nach der Referenz setzen

% Defines a \FloatBarrier command, beyond which floats may not
% pass; useful, for example, to ensure all floats for a section
% appear before the next \section command.
\usepackage[
	section		% "\section" command will be redefined with "\FloatBarrier"
]{placeins}
%
%%% Doc: ftp://tug.ctan.org/pub/tex-archive/macros/latex/contrib/subfig/subfig.pdf
% Incompatible: loads package capt-of. Loading of 'capt-of' afterwards will fail therefor
\usepackage{subfig} % Layout wird weiter unten festgelegt !

%%% Bilder von Text Umfliessen lassen : (empfehle wrapfig)
%
%%% Doc: ftp://tug.ctan.org/pub/tex-archive/macros/latex/contrib/wrapfig/wrapfig.sty
\usepackage{wrapfig}	        % defines wrapfigure and wrapfloat
%\setlength{\wrapoverhang}{\marginparwidth} % aeerlapp des Bildes ...
%\addtolength{\wrapoverhang}{\marginparsep} % ... in den margin
\setlength{\intextsep}{0.5\baselineskip} % Platz ober- und unterhalb des Bildes
% \intextsep ignoiert bei draft ???
%\setlength{\columnsep}{1em} % Abstand zum Text

%%% Doc: Documentation inside dtx Package
%\usepackage{floatflt}   	  % LaTeX2e Paket von 1996
                             % [rflt] - Standard float auf der rechten Seite

%%% Doc: ftp://tug.ctan.org/pub/tex-archive/macros/latex209/contrib/picins/picins.doc
%\usepackage{picins}          % LaTeX 2.09 Paket von 1992. aber Layout kombatibel


% Make float placement easier
\renewcommand{\floatpagefraction}{.75} % vorher: .5
\renewcommand{\textfraction}{.1}       % vorher: .2
\renewcommand{\topfraction}{.8}        % vorher: .7
\renewcommand{\bottomfraction}{.5}     % vorher: .3
\setcounter{topnumber}{3}              % vorher: 2
\setcounter{bottomnumber}{2}           % vorher: 1
\setcounter{totalnumber}{5}            % vorher: 3


%%% Doc: ftp://tug.ctan.org/pub/tex-archive/macros/latex/contrib/psfrag/pfgguide.pdf
% \usepackage{psfrag}	% Ersetzen von Zeichen in eps Bildern


%%% Doc: http://www.ctan.org/tex-archive/macros/latex/contrib/sidecap/sidecap.pdf
\usepackage[%
%	outercaption,%	(default) caption is placed always on the outside side
%	innercaption,% caption placed on the inner side
%	leftcaption,%  caption placed on the left side
	rightcaption,% caption placed on the right side
%	wide,%			caption of float my extend into the margin if necessary
%	margincaption,% caption set into margin
	ragged,% caption is set ragged
]{sidecap}

\renewcommand\sidecaptionsep{2em}
%\renewcommand\sidecaptionrelwidth{20}
\sidecaptionvpos{table}{c}
\sidecaptionvpos{figure}{c}


%% Diagramme mit LaTeX ===================================================
%

%%% Doc: ftp://tug.ctan.org/pub/tex-archive/macros/latex/contrib/pict2e/pict2e.pdf
% Neuimplementation der Picture Umgebung.
%
% The new package extends the existing LaTeX picture environment, using
% the familiar technique (cf. the graphics and color packages) of driver
% files.  The package documentation (pict2e.dtx) has a fair number of
% examples of use, showing where things are improved by comparison with
% the LaTeX picture environment.
% \usepackage{pict2e}

%%% Doc: ftp://tug.ctan.org/pub/tex-archive/macros/latex/contrib/curve2e/curve2e.pdf
% Extensions for package pict2e.
%\usepackage{curve2e}
%


% ~~~~~~~~~~~~~~~~~~~~~~~~~~~~~~~~~~~~~~~~~~~~~~~~~~~~~~~~~~~~~~~~~~~~~~~~
% misc packages
% ~~~~~~~~~~~~~~~~~~~~~~~~~~~~~~~~~~~~~~~~~~~~~~~~~~~~~~~~~~~~~~~~~~~~~~~~

\usepackage{makeidx}		% Index
\IfDraft{
  \usepackage{showidx}    % Indizierte Begriffe am Rand (Korrekturlesen)
}

%%% Doc: ftp://tug.ctan.org/pub/tex-archive/macros/latex/contrib/isodate/README
%%% Incompatible: draftcopy
% Tune the output format of dates.
%\usepackage{isodate}

%%% Doc: ftp://tug.ctan.org/pub/tex-archive/macros/latex/contrib/numprint/numprint.pdf
% Modify printing of numbers
%\usepackage{numprint}

%%% Doc: ftp://tug.ctan.org/pub/tex-archive/macros/latex/contrib/nomencl/nomencl.pdf
\usepackage[%
	german,
	%english
]{nomencl}[2005/09/22]

\usepackage[
	footnote,	% Full names appear in the footnote
	smaller,		% Print acronym in smaller fontsize
	printonlyused %
]{acronym}

% ~~~~~~~~~~~~~~~~~~~~~~~~~~~~~~~~~~~~~~~~~~~~~~~~~~~~~~~~~~~~~~~~~~~~~~~~
% verbatim packages
% ~~~~~~~~~~~~~~~~~~~~~~~~~~~~~~~~~~~~~~~~~~~~~~~~~~~~~~~~~~~~~~~~~~~~~~~~

%%% Doc: ftp://tug.ctan.org/pub/tex-archive/macros/latex/contrib/upquote/upquote.sty
\usepackage{upquote} % Setzt "richtige" Quotes in verbatim-Umgebung

%%% Doc: No Documentation
% \usepackage{verbatim} %Reimplemntation of the original verbatim

%%% Doc: http://www.cs.brown.edu/system/software/latex/doc/fancyvrb.pdf
% \usepackage{fancyvrb} % Superior Verbatim Class

%% Listings Paket ------------------------------------------------------
%%% Doc: ftp://tug.ctan.org/pub/tex-archive/macros/latex/contrib/listings/listings-1.3.pdf
% \usepackage{listings}
% \lstset{
%         basicstyle=\small\ttfamily, % Standardschrift
%         numbers=left,               % Ort der Zeilennummern
%         numberstyle=\tiny,          % Stil der Zeilennummern
%         stepnumber=2,               % Abstand zwischen den Zeilennummern
%         numbersep=5pt,              % Abstand der Nummern zum Text
%         tabsize=2,                  % Groesse von Tabs
%         extendedchars=true,         %
%         breaklines=true,            % Zeilen werden Umgebrochen
% %        keywordstyle=[1]\textbf,    % Stil der Keywords
% %        keywordstyle=[2]\textbf,    %
% %        keywordstyle=[3]\textbf,    %
% %        keywordstyle=[4]\textbf,    %
%         stringstyle=\color{stringcolor}, % Farbe der String
%         showspaces=false,           % Leerzeichen anzeigen ?
%         showtabs=false,             % Tabs anzeigen ?
%         showstringspaces=false      % Leerzeichen in Strings anzeigen ?
% }
% \lstloadlanguages{% Check Dokumentation for further languages ...
%         [Visual]Basic
%         %Pascal
%         %C
%         %C++
%         %XML
%         %HTML
% }

%%% Doc: ftp://tug.ctan.org/pub/tex-archive/macros/latex/contrib/examplep/eurotex_2005_examplep.pdf
% LaTeX Code und Ergebnis nebeneinander darstellen
%\usepackage{examplep}

% ~~~~~~~~~~~~~~~~~~~~~~~~~~~~~~~~~~~~~~~~~~~~~~~~~~~~~~~~~~~~~~~~~~~~~~~~
% science packages
% ~~~~~~~~~~~~~~~~~~~~~~~~~~~~~~~~~~~~~~~~~~~~~~~~~~~~~~~~~~~~~~~~~~~~~~~~

\usepackage{units}


% ~~~~~~~~~~~~~~~~~~~~~~~~~~~~~~~~~~~~~~~~~~~~~~~~~~~~~~~~~~~~~~~~~~~~~~~~
% fancy packages
% ~~~~~~~~~~~~~~~~~~~~~~~~~~~~~~~~~~~~~~~~~~~~~~~~~~~~~~~~~~~~~~~~~~~~~~~~

%%% Doc: No documentation - documented in 'The LaTeX Companion'
% \usepackage{fancybox}   % for shadowbox, ovalbox

%%% Doc: ftp://tug.ctan.org/pub/tex-archive/macros/latex/contrib/misc/framed.sty
% \usepackage{framed}
% \renewcommand\FrameCommand{\fcolorbox{black}{shadecolor}}

\makeatletter
\IfPackageLoaded{framed}{%
   \IfPackageLoaded{marginnote}{%
������\begingroup
         \g@addto@macro\framed{%
���������   \let\marginnoteleftadjust\FrameSep
������������   \let\marginnoterightadjust\FrameSep
         }
      \makeatother
� }
}
\makeatother



%%% Doc: No documentation - documented in 'The LaTeX Companion'
% \usepackage{boxedminipage}

%%% Doc: ftp://tug.ctan.org/pub/tex-archive/macros/latex/contrib/lettrine/doc/lettrine.pdf
% Dropping capitals
% \usepackage{lettrine}


% ~~~~~~~~~~~~~~~~~~~~~~~~~~~~~~~~~~~~~~~~~~~~~~~~~~~~~~~~~~~~~~~~~~~~~~~~
% layout packages
% ~~~~~~~~~~~~~~~~~~~~~~~~~~~~~~~~~~~~~~~~~~~~~~~~~~~~~~~~~~~~~~~~~~~~~~~~

%%% Diverse Pakete und Einstellungen =====================================

%%% Doc: Documentation inside dtx file
% Mehere Text-Spalten
\usepackage{multicol}

%\nonfrenchspacing     % liefert extra Platz hinter Satzenden.
                       % Fuer deutschen Text standardmaessig ausgeschaltet!


\usepackage{ellipsis}  % >>Intelligente<< \dots

%% Zeilenabstand =========================================================
%
%%% Doc: ftp://tug.ctan.org/pub/tex-archive/macros/latex/contrib/setspace/setspace.sty
\usepackage{setspace}
%\onehalfspacing		% 1,5-facher Abstand
%\doublespacing		% 2-facher Abstand
% hereafter load 'typearea' again

%% Seitenlayout ==========================================================
%
% Layout laden um im Dokument den Befehl \layout nutzen zu koennen
%%% Doc: no documentation
%\usepackage[verbose]{layout}
%

% Layout mit 'geometry'
%%% Doc: ftp://tug.ctan.org/pub/tex-archive/macros/latex/contrib/geometry/manual.pdf
% \usepackage{geometry}

\IfPackageLoaded{geometry}{%
\geometry{%
%%% Paper Groesse
   a4paper, % Andere a0paper, a1paper, a2paper, a3paper, , a5paper, a6paper,
            % b0paper, b1paper, b2paper, b3paper, b4paper, b5paper, b6paper
            % letterpaper, executivepaper, legalpaper
   %screen,  % a special paper size with (W,H) = (225mm,180mm)
   %paperwidth=,
   %paperheight=,
   %papersize=, %{ width , height }
   %landscape,  % Querformat
   portrait,    % Hochformat
%%% Koerper Groesse
   %hscale=,      % ratio of width of total body to \paperwidth
                  % hscale=0.8 is equivalent to width=0.8\paperwidth. (0.7 by default)
   %vscale=,      % ratio of height of total body to \paperheight
                  % vscale=0.9 is equivalent to height=0.9\paperheight.
   %scale=,       % ratio of total body to the paper. scale={ h-scale , v-scale }
   %totalwidth=,    % width of total body % (Generally, width >= textwidth)
   %totalheight=,   % height of total body, excluding header and footer by default
   %total=,        % total={ width , height }
   %textwidth=,    % modifies \textwidth, the width of body
   %textheight=,   % modifies \textheight, the height of body
   %body=,        % { width , height } sets both \textwidth and \textheight of the body of page.
   lines=45,       % enables users to specify \textheight by the number of lines.
   %includehead,  % includes the head of the page, \headheight and \headsep, into total body.
   %includefoot,  % includes the foot of the page, \footskip, into body.
   %includeheadfoot, % sets both includehead and includefoot to true
   %includemp,    % includes the margin notes, \marginparwidth and \marginparsep, into body
   %includeall,   % sets both includeheadfoot and includemp to true.
   %ignorehead,   % disregards the head of the page, headheight and headsep in determining vertical layout
   %ignorefoot,   % disregards the foot of page, footskip, in determining vertical layout
   %ignoreheadfoot, % sets both ignorehead and ignorefoot to true.
   %ignoremp,     % disregards the marginal notes in determining the horizontal margins
   ignoreall,     % sets both ignoreheadfoot and ignoremp to true
   heightrounded, % This option rounds \textheight to n-times (n: an integer) of \baselineskip
   %hdivide=,     % { left margin , width , right margin }
                  % Note that you should not specify all of the three parameters
   %vdivide=,     % { top margin , height , bottom margin }
   %divide=,      % ={A,B,C} %  is interpreted as hdivide={A,B,C} and vdivide={A,B,C}.
%%% Margin
   %left=,        % left margin (for oneside) or inner margin (for twoside) of total body
                  % alias: lmargin, inner
   %right=,       % right or outer margin of total body
                  % alias: rmargin outer
   %top=,         % top margin of the page.
                  % Alias : tmargin
   %bottom=,      % bottom margin of the page
                  % Alias : bmargin
   %hmargin=,     % left and right margin. hmargin={ left margin , right margin }
   %vmargin=,     % top and bottom margin. vmargin={ top margin , bottom margin }
   %margin=,      % margin={A,B} is equivalent to hmargin={A,B} and vmargin={A,B}
   %hmarginratio, % horizontal margin ratio of left (inner) to right (outer).
   %vmarginratio, % vertical margin ratio of top to bottom.
   %marginratio,  % marginratio={ horizontal ratio , vertical ratio }
   %hcentering,   % sets auto-centering horizontally and is equivalent to hmarginratio=1:1
   %vcentering,   % sets auto-centering vertically and is equivalent to vmarginratio=1:1
   %centering,    % sets auto-centering and is equivalent to marginratio=1:1
   twoside,       % switches on twoside mode with left and right margins swapped on verso pages.
   %asymmetric,   % implements a twosided layout in which margins are not swapped on alternate pages
                  % and in which the marginal notes stay always on the same side.
   bindingoffset=5mm,  % removes a specified space for binding
%%% Dimensionen
   %headheight=,  % Alias:  head
   %headsep=,     % separation between header and text
   %footskip=,    % distance separation between baseline of last line of text and baseline of footer
                  % Alias: foot
   %nohead,       % eliminates spaces for the head of the page
                  % equivalent to both \headheight=0pt and \headsep=0pt.
   %nofoot,       % eliminates spaces for the foot of the page
                  % equivalent to \footskip=0pt.
   %noheadfoot,   % equivalent to nohead and nofoot.
   %footnotesep=, % changes the dimension \skip\footins,.
                  % separation between the bottom of text body and the top of footnote text
   marginparwidth=0pt, % width of the marginal notes
                  % Alias: marginpar
   %marginparsep=,% separation between body and marginal notes.
   %nomarginpar,  % shrinks spaces for marginal notes to 0pt
   %columnsep=,   % the separation between two columns in twocolumn mode.
   %hoffset=,
   %voffset=,
   %offset=,      % horizontal and vertical offset.
                  % offset={ hoffset , voffset }
   %twocolumn,    % twocolumn=false denotes onecolumn
   twoside,
   textwidth=400pt,   % sets \textwidth directly
   %textheight=,  % sets \textheight directly
   %reversemp,    % makes the marginal notes appear in the left (inner) margin
                  % Alias: reversemarginpar
}
} % Endif

% - Anzeigen des Layouts -
\IfPackageLoaded{geometry}{%
   %\geometry{showframe}
}

% Layout mit 'typearea'
%%% Doc: ftp://tug.ctan.org/pub/tex-archive/macros/latex/contrib/koma-script/scrguide.pdf

\IfPackageLoaded{typearea}{% Wenn typearea geladen ist
   \IfPackageNotLoaded{geometry}{% aber nicht geometry
      \typearea[current]{last}
   }
}

% BCOR
%    current  % Satzspiegelberechnung mit dem aktuell g�ltigen BCOR-Wert erneut
%             % durchf�hren.
% DIV
%    calc     % Satzspiegelberechnung einschlie�lich Ermittlung eines guten
%             % DIV-Wertes erneut durchf�hren.
%    classic  % Satzspiegelberechnung nach dem
%             % mittelalterlichen Buchseitenkanon
%             % (Kreisberechnung) erneut durchf�hren.
%    current  % Satzspiegelberechnung mit dem aktuell g�ltigen DIV-Wert erneut
%             % durchf�hren.
%    default  % Satzspiegelberechnung mit dem Standardwert f�r das aktuelle
%             % Seitenformat und die aktuelle Schriftgr��e erneut durchf�hren.
%             % Falls kein Standardwert existiert calc anwenden.
%    last     % Satzspiegelberechnung mit demselben DIV -Argument, das beim
%             % letzten Aufruf angegeben wurde, erneut durchf�hren


%\usepackage[colorgrid,texcoord,gridunit=mm]{showframe}

\raggedbottom     % Variable Seitenhoehen zulassen

% Farben ================================================================

\IfDefined{definecolor}{%

% Farbe der Ueberschriften
%\definecolor{sectioncolor}{RGB}{0, 51, 153} % Blau
%\definecolor{sectioncolor}{RGB}{0, 25, 152}    % Blau (dunkler))
\definecolor{sectioncolor}{RGB}{0, 0, 0}    % Schwarz
%
% Farbe des Textes
\definecolor{textcolor}{RGB}{0, 0, 0}        % Schwarz
%
% Farbe fuer grau hinterlegte Boxen (fuer Paket framed.sty)
\definecolor{shadecolor}{gray}{0.90}

% Farben fuer die Links im PDF
\definecolor{pdfurlcolor}{rgb}{0,0,0.6}
\definecolor{pdffilecolor}{rgb}{0.7,0,0}
\definecolor{pdflinkcolor}{rgb}{0,0,0.6}
\definecolor{pdfcitecolor}{rgb}{0,0,0.6}

%% PDF-Linkfarben auf schwarz f�r den Druck:
% \definecolor{pdfurlcolor}{rgb}{0,0,0}
% \definecolor{pdffilecolor}{rgb}{0,0,0}
% \definecolor{pdflinkcolor}{rgb}{0,0,0}
% \definecolor{pdfcitecolor}{rgb}{0,0,0}


% Farben fuer Listings
\colorlet{stringcolor}{green!40!black!100}
\colorlet{commencolor}{blue!0!black!100}

} % Endif

%% Aussehen der URLS======================================================

%fuer URL (nur wenn url geladen ist)
\IfDefined{urlstyle}{
	\urlstyle{tt} %sf
}

%% Kopf und Fusszeilen====================================================
%%% Doc: ftp://tug.ctan.org/pub/tex-archive/macros/latex/contrib/koma-script/scrguide.pdf

\usepackage[%
   % headtopline,
   % plainheadtopline,
   % headsepline,
   % plainheadsepline,
   % footsepline,
   % plainfootsepline,
   % footbotline,
   % plainfootbotline,
   % ilines,
   % clines,
   % olines,
   automark,
   % autooneside,% ignore optional argument in automark at oneside
   komastyle,
   % standardstyle,
   % markuppercase,
   % markusedcase,
   nouppercase,
]{scrpage2}

%\usepackage[%
%   automark,         % automatische Aktualisierung der Kolumnentitel
%   nouppercase,      % Grossbuchstaben verhindern
%   %markuppercase    % Grossbuchstaben erzwingen
%   %markusedcase     % vordefinierten Stil beibehalten
%   %komastyle,       % Stil von Koma Script
%   %standardstyle,   % Stil der Standardklassen
%]{scrpage2}

\IfElseChapterDefined{%
   \pagestyle{scrheadings} % Seite mit Headern
}{
   \pagestyle{scrplain} % Seiten ohne Header
}
%\pagestyle{empty} % Seiten ohne Header
%
% loescht voreingestellte Stile
\clearscrheadings
\clearscrplain
%
% Was steht wo...
\IfElseChapterDefined{
   % Oben aussen: Kapitel und Section
   % Unten aussen: Seitenzahl
   % \ohead{\headmark} % Oben au�en: Setzt Kapitel und Section automatisch
   % \ofoot[\pagemark]{\pagemark}
   % oder...
   % Oben aussen: Seitenzahlen
   % Oben innen: Kapitel und Section
   \ohead{\pagemark}
   \ihead{\headmark}
   \ofoot[\pagemark]{} % Au�en unten: Seitenzahlen bei plain
}{
   \cfoot[\pagemark]{\pagemark} % Mitte unten: Seitenzahlen bei plain
}
% Vollstaendige Liste der moeglichen Positionierungen
% \lehead[scrplain-links-gerade]{scrheadings-links-gerade}
% \cehead[scrplain-mittig-gerade]{scrheadings-mittig-gerade}
% \rehead[scrplain-rechts-gerade]{scrheadings-rechts-gerade}
% \lefoot[scrplain-links-gerade]{scrheadings-links-gerade}
% \cefoot[scrplain-mittig-gerade]{scrheadings-mittig-gerade}
% \refoot[scrplain-rechts-gerade]{scrheadings-rechts-gerade}
% \lohead[scrplain-links-ungerade]{scrheadings-links-ungerade}
% \cohead[scrplain-mittig-ungerade]{scrheadings-mittig-ungerade}
% \rohead[scrplain-rechts-ungerade]{scrheadings-rechts-ungerade}
% \lofoot[scrplain-links-ungerade]{scrheadings-links-ungerade}
% \cofoot[scrplain-mittig-ungerade]{scrheadings-mittig-ungerade}
% \rofoot[scrplain-rechts-ungerade]{scrheadings-rechts-ungerade}
% \ihead[scrplain-innen]{scrheadings-innen}
% \chead[scrplain-zentriert]{scrheadings-zentriert}
% \ohead[scrplain-au�en]{scrheadings-au�en}
% \ifoot[scrplain-innen]{scrheadings-innen}
% \cfoot[scrplain-zentriert]{scrheadings-zentriert}
% \ofoot[scrplain-au�en]{scrheadings-au�en}


%\usepackage{lastpage} % Stellt 'LastPage' zur Verfuegung
%\cfoot[Seite \pagemark~von \pageref{LastPage}]{} % Seitenzahl von Anzahl Seiten

% Angezeigte Abschnitte im Header
\IfElseChapterDefined{
   \automark[section]{chapter} %[rechts]{links}
}{
   \automark[subsection]{section} %[rechts]{links}
}
%
% Linien (moegliche Kombination mit Breiten)
\IfChapterDefined{
   %\setheadtopline{}     % modifiziert die Parameter fuer die Linie ueber dem
   							  %	 Seitenkopf
   \setheadsepline{.4pt}[\color{black}]
                         % modifiziert die Parameter fuer die Linie zwischen
                         % Kopf und Textk�rper
   %\setfootsepline{}    % modifiziert die Parameter fuer die Linie zwischen
   							 % Text und Fu�
   %\setfootbotline{}    % modifiziert die Parameter fuer die Linie unter dem
   							 % Seitenfuss
}



% Groesse des Headers

% Breite von Kopf und Fusszeile einstellen
% \setheadwidth[Verschiebung]{Breite}
% \setfootwidth[Verschiebung]{Breite}
% m�gliche Werte
% paper - die Breite des Papiers
% page - die Breite der Seite
% text - die Breite des Textbereichs
% textwithmarginpar - die Breite des Textbereichs inklusive dem Seitenrand
% head - die aktuelle Breite des Seitenkopfes
% foot - die aktuelle Breite des Seitenfusses
\setheadwidth[0pt]{text}
\setfootwidth[0pt]{text}

%\setheadwidth[]%
%{%
%   paper % width of paper
%   page  % width of page (paper - BCOR)
%   text  % \textwidth
%   textwithmarginpar % width of text plus margin
%   head  % current width of head
%   foot  % current width of foot
%}%


%% Fussnoten =============================================================
% Keine hochgestellten Ziffern in der Fussnote (KOMA-Script-spezifisch):
\deffootnote{1.5em}{1em}{\makebox[1.5em][l]{\thefootnotemark}}
\addtolength{\skip\footins}{\baselineskip} % Abstand Text <-> Fussnote

\setlength{\dimen\footins}{10\baselineskip} % Beschraenkt den Platz von Fussnoten auf 10 Zeilen

\interfootnotelinepenalty=10000 % Verhindert das Fortsetzen von
                                % Fussnoten auf der gegen�berligenden Seite


%% Schriften (Sections )==================================================

\IfElsePackageLoaded{fourier}{
   \newcommand\SectionFontStyle{\rmfamily}
}{
   \newcommand\SectionFontStyle{\sffamily}
}

% -- Koma Schriften --
\IfChapterDefined{%
   \setkomafont{chapter}{\huge\SectionFontStyle}    % Chapter
}

\setkomafont{sectioning}{\SectionFontStyle}
%\setkomafont{part}{\usekomafont{sectioning}}
%\setkomafont{section}{\usekomafont{sectioning}}
%\setkomafont{subsection}{\usekomafont{sectioning}}
%\setkomafont{subsubsection}{\usekomafont{sectioning}}
%\setkomafont{paragraph}{\usekomafont{sectioning}}
%\setkomafont{subparagraph}{\usekomafont{sectioning}}

\setkomafont{descriptionlabel}{\itshape}

%\setkomafont{caption}{\normalfont}
%\setkomafont{captionlabel}{\normalfont}

%\setkomafont{dictum}{}
%\setkomafont{dictumauthor}{}
%\setkomafont{dictumtext}{}
%\setkomafont{disposition}{}
%\setkomafont{footnote}{}
%\setkomafont{footnotelabel}{}
%\setkomafont{footnotereference}{}
%\setkomafont{minisec}{}

%\setkomafont{partnumber}{\bfseries\SectionFontStyle}
%\setkomafont{partentrynumber}{}
%\setkomafont{chapterentrypagenumber}{}
%\setkomafont{sectionentrypagenumber}{}

\setkomafont{pageheadfoot}{\normalfont\normalcolor\small\sffamily}
\setkomafont{pagenumber}{\bfseries\usekomafont{sectioning}}

%\setkomafont{partentry}{\usekomafont{sectioning}\large}
%\setkomafont{chapterentry}{\usekomafont{sectioning}}
%\setkomafont{sectionentry}{\usekomafont{sectioning}}

%%% --- Titlepage ---
%\setkomafont{subject}{}
%\setkomafont{subtitle}{}
%\setkomafont{title}{}



\addtokomafont{sectioning}{\color{sectioncolor}} % Farbe der Ueberschriften
\IfChapterDefined{%
	\addtokomafont{chapter}{\color{sectioncolor}} % Farbe der Ueberschriften
}
\renewcommand*{\raggedsection}{\raggedright} % Titelzeile linksbuendig, haengend
%
%% UeberSchriften (Chapter und Sections) =================================

%%% Remove Space above Chapter.
%%% (NOT recommanded!)
%% Space above Chapter Title
% \renewcommand*{\chapterheadstartvskip}{\vspace{1\baselineskip}}%
%% Space below Chapter Title
% \renewcommand*{\chapterheadendvskip}{\vspace{0.5\baselineskip}}%

% -- Ueberschriften komlett Umdefinieren --
%%% Doc: ftp://tug.ctan.org/pub/tex-archive/macros/latex/contrib/titlesec/titlesec.pdf
%\usepackage{titlesec}

% -- Section Aussehen veraendern --
% --------------------------------
%% -> Section mit Unterstrich
% \titleformat{\section}
%   [hang]%[frame]display
%   {\usekomafont{sectioning}\Large}
%  {\thesection}
%   {6pt}
%   {}
%   [\titlerule \vspace{0.5\baselineskip}]
% --------------------------------

% -- Chapter Aussehen veraendern --
% --------------------------------
%--> Box mit (Kapitel + Nummer ) +  Name
% \titleformat{\chapter}[display]     % {command}[shape]
%   {\usekomafont{chapter}\filcenter} % format
%   {                                 % label
%   {\fcolorbox{black}{shadecolor}{
%   {\huge\chaptertitlename\mbox{\hspace{1mm}}\thechapter}
%   }}}
%   {1pc}                             % sep (from chapternumber)
%   {\vspace{1pc}}                    % {before}[after] (before chaptertitle and after)
% --------------------------------
%--> Kapitel + Nummer + Trennlinie + Name + Trennlinie
%\titleformat{\chapter}[display]	% {command}[shape]
%  {\usekomafont{chapter}\Large \color{black}}	% format
%  {   										% label
%  \LARGE\MakeUppercase{\chaptertitlename} \Huge \thechapter \filright%
%  }%}
%  {1pt}										% sep (from chapternumber)
%  {\titlerule \vspace{0.9pc} \filright \color{sectioncolor}}   % {before}[after] (before chaptertitle and after)
%  [\color{black} \vspace{0.9pc} \filright {\titlerule}]


%%% Doc: No documentation
% Indent first paragraph after section header
% \usepackage{indentfirst}

%% Captions (Schrift, Aussehen) ==========================================

% % Folgende Befehle werden durch das Paket caption und subfig ersetzt !
% \setcapindent{1em} % Einrueckung der Beschriftung
% \setkomafont{caption}{\color{black}\small\sffamily\RaggedRight}  % Schrift fuer Caption
% \setkomafont{captionlabel}{\color{black}\small}   % Schrift fuer 'Abbildung' usw.

%%% Doc: ftp://tug.ctan.org/pub/tex-archive/macros/latex/contrib/caption/caption.pdf
\usepackage{caption}
% Aussehen der Captions
\captionsetup{
   margin = 10pt,
   font = {small,rm},
   labelfont = {small,bf},
   format = plain, % oder 'hang'
   indention = 0em,  % Einruecken der Beschriftung
   labelsep = colon, %period, space, quad, newline
   justification = RaggedRight, % justified, centering
   singlelinecheck = true, % false (true=bei einer Zeile immer zentrieren)
   position = bottom %top
}
%%% Bugfix Workaround
\DeclareCaptionOption{parskip}[]{}
\DeclareCaptionOption{parindent}[]{}

% Aussehen der Captions fuer subfigures (subfig-Paket)
\IfPackageLoaded{subfig}{
 \captionsetup[subfloat]{%
   margin = 10pt,
   font = {small,rm},
   labelfont = {small,bf},
   format = plain, % oder 'hang'
   indention = 0em,  % Einruecken der Beschriftung
   labelsep = space, %period, space, quad, newline
   justification = RaggedRight, % justified, centering
   singlelinecheck = true, % false (true=bei einer Zeile immer zentrieren)
   position = bottom, %top
   labelformat = parens % simple, empty % Wie die Bezeichnung gesetzt wird
 }
}

% Aendern der Bezeichnung fuer Abbildung und Tabelle
% \addto\captionsngerman{% "captionsgerman" fuer alte  Rechschreibung
%   \renewcommand{\figurename}{Abb.}%
%   \renewcommand{\tablename}{Tab.}%
% }

% Caption fuer nicht fliessende Umgebungen
%%% Doc: ftp://tug.ctan.org/pub/tex-archive/macros/latex/contrib/misc/capt-of.sty
\IfPackageNotLoaded{caption}{
	\usepackage{capt-of} % only load when caption is not loaded. Otherwise compiling will fail.
	%Usage: \captionof{table}[short Titel]{long Titel}
}
%


%%% Doc: ftp://tug.ctan.org/pub/tex-archive/macros/latex/contrib/mcaption/mcaption.pdf
% Captions in Margins
% \usepackage[
% 	top,
% 	bottom
% ]{mcaption}

%%% Example:
% \begin{figure}
%   \begin{margincap}[short caption]{margin caption}
%     \centering
%     \includegraphics{picture}
%   \end{margincap}
% \end{figure}



% \numberwithin{figure}{chapter} %Befehl zum Kapitelweise Nummerieren der Bilder, setzt `amsmath' vorraus
% \numberwithin{table}{chapter}  %Befehl zum Kapitelweise Nummerieren der Tabellen, setzt `amsmath' vorraus

%% Inhaltsverzeichnis (Schrift, Aussehen) sowie weitere Verzeichnisse ====

\setcounter{secnumdepth}{2}    % Abbildungsnummerierung mit groesserer Tiefe
\setcounter{tocdepth}{2}		 % Inhaltsverzeichnis mit groesserer Tiefe
%

% Inhalte von List of Figures
\IfPackageLoaded{subfig}{
	\setcounter{lofdepth}{1}  %1 = nur figures, 2 = figures + subfigures
}

% -------------------------------------------------------

% Aussehen des Inhaltsverzeichnisses: tocloft
%%% Doc: ftp://tug.ctan.org/pub/tex-archive/macros/latex/contrib/tocloft/tocloft.pdf
%% Laden mit Option subfigure in Abhaengigkeit vom Paket subfigure und subfig
% \IfElsePackageLoaded{subfig}
% 	% IF subfig
% 	{\usepackage[subfigure]{tocloft}}{
% 	% ELSE
% 	\IfElsePackageLoaded{subfigure}
% 		% IF subfigure
% 		{\usepackage[subfigure]{tocloft}}
% 	   % Else (No subfig nor subfigure)
% 		{\usepackage{tocloft}}
% 	}
%
% %TOCLOFT zerstoert Layout der Ueberschriften von TOC, LOT, LOF
% \IfPackageLoaded{tocloft}{
% %
% %%%% Layout Matthias Pospiech (alles serifenlos)
% \IfChapterDefined{%
% 	\renewcommand{\cftchappagefont}{\bfseries\sffamily}  % Kapitel Seiten Schrift
% 	\renewcommand{\cftchapfont}{\bfseries\sffamily}      % Kapitel Schrift
% }
% \renewcommand{\cftsecpagefont}{\sffamily}            % Section Seiten Schrift
% \renewcommand{\cftsubsecpagefont}{\sffamily}         % Subsectin Seiten Schrift
% \renewcommand{\cftsecfont}{\sffamily}                % Section Schrift
% \renewcommand{\cftsubsecfont}{\sffamily}             % Subsection Schrift
%
% %%%% Layout aus Typokurz:
% % % Seitenzahlen direkt hinter TOC-Eintrag:
% % % Ebene \chapter
% % \renewcommand{\cftchapleader}{}
% % \renewcommand{\cftchapafterpnum}{\cftparfillskip}
% % % Ebene \section
% % \renewcommand{\cftsecleader}{}
% % \renewcommand{\cftsecafterpnum}{\cftparfillskip}
% % % Ebene \subsection
% % \renewcommand{\cftsubsecleader}{}
% % \renewcommand{\cftsubsecafterpnum}{\cftparfillskip}
% % % Abstaende vor Eintraegen im TOC verkleinern
% % \setlength{\cftbeforesecskip}{.4\baselineskip}
% % \setlength{\cftbeforesubsecskip}{.1\baselineskip}
% }
% % Ende tocloft Einstellungen --------------

%%% Doc: ftp://tug.ctan.org/pub/tex-archive/macros/latex/contrib/ms/multitoc.dvi
% TOC in mehreren Spalten setzen
%\usepackage[toc]{multitoc}

% -------------------------

%%Schriften fuer Minitoc (Inhaltsverzeichnis vor jedem Kapitel)
%%% Doc: ftp://tug.ctan.org/pub/tex-archive/macros/latex/contrib/minitoc/minitoc.pdf
%\IfElseChapterDefined{%
% \usepackage{minitoc}
% \setlength{\mtcindent}{0em} % default: 24pt
% \setcounter{minitocdepth}{2}
% \setlength{\mtcskipamount}{\bigskipamount}
% \mtcsettitlefont{minitoc}{\normalsize\SectionFontStyle}
% \mtcsetfont{minitoc}{*}{\small\SectionFontStyle} %\color{textcolor}
% \mtcsetfont{minitoc}{section}{\small\SectionFontStyle}
% \mtcsetfont{minitoc}{subsection}{\small\SectionFontStyle}
% \mtcsetfont{minitoc}{subsubsection}{\small\SectionFontStyle}
%}{
%�\usepackage{minitoc}
%�\setlength{\stcindent}{0pt} %default
%�\setcounter{secttocdepth}{2} %default
%�\mtcsettitlefont{secttoc}{\SectionFontStyle}
%�\mtcsetfont{secttoc}{*}{\small\SectionFontStyle}%
%�\mtcsetfont{secttoc}{subsection}{\small\SectionFontStyle}
%�\mtcsetfont{secttoc}{subsubsection}{\small\SectionFontStyle}
%}

% Packages that MUST be loaded before minitoc !
% hyperref, caption, sectsty, varsects, fncychap, hangcaption, quotchap, romannum, sfheaders, alnumsec, captcont


%% Index & Co. ===========================================================
% gibts dafuer noch eine sauberere Loesung ?
%%%%%%%% Index zweispaltig %%%%%%%
% \makeatletter
% \renewenvironment{theindex}{%
% \setlength{\columnsep}{2em}
% \begin{multicols}{2}[\section*{\indexname}]
% \parindent\z@
% \parskip\z@ \@plus .3\p@\relax
% \let\item\@idxitem}%
% {\end{multicols}\clearpage}
% \makeatother
%%%%%%%%%%%%%%%%%%%%%%%%%%%%%%%%%%


% ~~~~~~~~~~~~~~~~~~~~~~~~~~~~~~~~~~~~~~~~~~~~~~~~~~~~~~~~~~~~~~~~~~~~~~~~
% PDF related packages
% ~~~~~~~~~~~~~~~~~~~~~~~~~~~~~~~~~~~~~~~~~~~~~~~~~~~~~~~~~~~~~~~~~~~~~~~~

%%% Doc: ftp://tug.ctan.org/pub/tex-archive/macros/latex/contrib/microtype/microtype.pdf
% Optischer Randausgleich mit pdfTeX
\ifpdf
\usepackage[%
	expansion=true, % better typography, but with much larger PDF file.
	protrusion=true
]{microtype}
\fi


%% Use only instead of hyperref !
% \usepackage[%
%    %ref,     % verweist auf Abschnitte
%    pageref, % verweist auf Seiten
% ]{backref} % Links in BiB back to Citation page/section (can be loaded by hyperref too)


%%% Doc: ftp://tug.ctan.org/pub/tex-archive/macros/latex/contrib/hyperref/doc/manual.pdf
%\usepackage[
%   % Farben fuer die Links
%   colorlinks=true,         % Links erhalten Farben statt Kaeten
%   urlcolor=pdfurlcolor,    % \href{...}{...} external (URL)
%   filecolor=pdffilecolor,  % \href{...} local file
%   linkcolor=pdflinkcolor,  %\ref{...} and \pageref{...}
%   citecolor=pdfcitecolor,  %
%   % Links
%   raiselinks=true,			 % calculate real height of the link
%   breaklinks,              % Links berstehen Zeilenumbruch
%   backref=page,            % Backlinks im Literaturverzeichnis (section, slide, page, none)
%   pagebackref=true,        % Backlinks im Literaturverzeichnis mit Seitenangabe
%   verbose,
%   hyperindex=true,         % backlinkex index
%   linktocpage=true,        % Inhaltsverzeichnis verlinkt Seiten
%   hyperfootnotes=false,     % Keine Links auf Fussnoten
%   % Bookmarks
%   bookmarks=true,          % Erzeugung von Bookmarks fuer PDF-Viewer
%   bookmarksopenlevel=1,    % Gliederungstiefe der Bookmarks
%   bookmarksopen=true,      % Expandierte Untermenues in Bookmarks
%   bookmarksnumbered=true,  % Nummerierung der Bookmarks
%   bookmarkstype=toc,       % Art der Verzeichnisses
%   % Anchors
%   plainpages=false,        % Anchors even on plain pages ?
%   pageanchor=true,         % Pages are linkable
%   % PDF Informationen
%   pdftitle={},             % Titel
%   pdfauthor={},            % Autor
%   pdfcreator={LaTeX, hyperref, KOMA-Script}, % Ersteller
%   %pdfproducer={pdfeTeX 1.10b-2.1} %Produzent
%   pdfdisplaydoctitle=true, % Dokumententitel statt Dateiname im Fenstertitel
%   pdfstartview=FitH,       % Dokument wird Fit Width geaefnet
%   pdfpagemode=UseOutlines, % Bookmarks im Viewer anzeigen
%   pdfpagelabels=true,           % set PDF page labels
%   pdfpagelayout=TwoPageRight, % zweiseitige Darstellung: ungerade Seiten
%   									 % rechts im PDF-Viewer
%   %pdfpagelayout=SinglePage, % einseitige Darstellung
%]{hyperref}

\IfPackageLoaded{backref}{
   % % Change Layout of Backref
   \renewcommand*{\backref}[1]{%
   	% default interface
   	% #1: backref list
   	%
   	% We want to use the alternative interface,
   	% therefore the definition is empty here.
   }%
   \renewcommand*{\backrefalt}[4]{%
   	% alternative interface
   	% #1: number of distinct back references
   	% #2: backref list with distinct entries
   	% #3: number of back references including duplicates
   	% #4: backref list including duplicates
   	\mbox{(Zitiert auf %
   	\ifnum#1=1 %
		   Seite~%
	   \else
   		Seiten~%
   	\fi
   	#2)}%
   }
}

%%% Doc: ftp://tug.ctan.org/pub/tex-archive/macros/latex/contrib/oberdiek/hypcap.pdf
% Links auf Gleitumgebungen springen nicht zur Beschriftung,
% sondern zum Anfang der Gleitumgebung
\IfPackageLoaded{hyperref}{%
	\usepackage[figure,table]{hypcap}
}

% Auch Abbildung und nicht nur die Nummer wird zum Link (abgeleitet
% aus Posting von Heiko Oberdiek (d09n5p$9md$1@news.BelWue.DE);
% Verwendung: In \abbvref{label} ist ein Beispiel dargestellt
\providecommand*{\figrefname}{Abbildung }
\newcommand*{\figref}[1]{%
  \hyperref[fig:#1]{\figrefname{}}\ref{fig:#1}%
}
% ebenso bei Tabellen
\providecommand*{\tabrefname}{Tabelle~}
\newcommand*{\tabref}[1]{%
  \hyperref[tab:#1]{\tabrefname{}}\ref{tab:#1}%
}
% und Abschnitten
\providecommand*{\secrefname}{Abschnitt }
\newcommand*{\secref}[1]{%
  \hyperref[sec:#1]{\secrefname{}}\ref{sec:#1}%
}
% und Kapiteln
\providecommand*{\chaprefname}{Kapitel~}
\newcommand*{\chapref}[1]{%
  \hyperref[chap:#1]{\chaprefname{}}\ref{chap:#1}%
}

%%% Doc: ftp://tug.ctan.org/pub/tex-archive/macros/latex/contrib/pdfpages/pdfpages.pdf
\usepackage{pdfpages} % Include pages from external PDF documents in LaTeX documents

%%% Doc: ftp://tug.ctan.org/pub/tex-archive/macros/latex/contrib/oberdiek/pdflscape.sty
%\usepackage{pdflscape} %  Querformat mit PDF
%
% Pakete Laden die nach Hyperref geladen werden sollen
\LoadPackagesNow % (ltxtable, tabularx)

% ~~~~~~~~~~~~~~~~~~~~~~~~~~~~~~~~~~~~~~~~~~~~~~~~~~~~~~~~~~~~~~~~~~~~~~~~
% Zus�tzliche Pakete
% ~~~~~~~~~~~~~~~~~~~~~~~~~~~~~~~~~~~~~~~~~~~~~~~~~~~~~~~~~~~~~~~~~~~~~~~~
%%% Doc: ftp://tug.ctan.org/pub/tex-archive/macros/latex/contrib/hyphenat/hyphenat.pdf
% According to documentation the font warnings can be ignored
%\usepackage[htt]{hyphenat} % enable hyphenation of typewriter text word (\textt).

%% Komprimierung von Bildern in PDF ausschalten
% \ifpdf
%    \pdfcompresslevel=0
% \fi
% ~~~~~~~~~~~~~~~~~~~~~~~~~~~~~~~~~~~~~~~~~~~~~~~~~~~~~~~~~~~~~~~~~~~~~~~~
% end of preambel
% ~~~~~~~~~~~~~~~~~~~~~~~~~~~~~~~~~~~~~~~~~~~~~~~~~~~~~~~~~~~~~~~~~~~~~~~~
%\IfPackageLoaded{fancyvrb}{
%	\DefineShortVerb{\|} % Nur mit fancyvrb zusammen laden!
%}


% Auszufuehrende Befehle  ------------------------------------------------
\IfDefined{makeindex}{\makeindex}
\IfDefined{makenomenclature}{\makenomenclature}
% \IfPackageLoaded{minitoc}{\IfElseUnDefined{chapter}{\dosecttoc}{\dominitoc}}
\IfPackageLoaded{minitoc}{\IfElseUnDefined{chapter}{\dosecttoc}{\dominitoc}}


\listfiles
%------------------------------------------------------------------------



%
%%%% Neue Befehle
% % redefine \textmu to other mu commands usefull inside text
% \renewcommand{\textmu}{$\upmu$}

% Caption with defined width
\newcommand{\wcaption}[2]{%
   \begin{minipage}{#1}%
   \caption{#2}%
   \end{minipage}%
}

\newcommand{\figureref}[1]{(Abbildung \ref{#1})}%
\newcommand{\eqnref}[1]{(\ref{#1})}%

% Command for margin text with usefull style
%\newcommand{\marginlabel}[1]{\mbox{}\marginline{\hspace{0pt}\footnotesize\sffamily #1}}%
\newcommand{\marginlabel}[1]{\marginnote{#1}}%

%\newcommand{\comment}[1]{\marginnote{#1}}%

% Enable space for figures that extent into the margin (right and/or leftside)
% Can be used inside a figure
% Note: sidecap defines a similar environment 'wide' !
\newenvironment{widespace}[2]{%
   \begin{list}{}{%
      \setlength{\topsep}{0pt}%
      \setlength{\leftmargin}{#1}%
      \setlength{\rightmargin}{#2}%
      \setlength{\listparindent}{\parindent}%
      \setlength{\itemindent}{\parskip}%
   }%
   \item[]%
}%
{%
   \end{list}%
}%

\newlength{\marginwidth}
\setlength{\marginwidth}{\marginparwidth}
\addtolength{\marginwidth}{\marginparsep}

%% Beispiel:
% \begin{figure}
% \begin{widespace}{-\marginwidth}{0pt}
%  \subfloat[Bergzebrastute]
%  {\includegraphics[width=0.45\linewidth]{../Bilder/Eingewoehnung2.jpg}}
%  \hspace*{1em}
%  \subfloat[Morro Moco]
%  {\includegraphics[width=0.45\linewidth]{../Bilder/bergzebra2.jpg}}
% \end{widespace}
% \end{figure}


% quantum optics - Latex Commands: Math **********************************
% ------------------------------------------------------------------------
% by: Matthias Pospiech
%%%%%%%%%%%%%%%%%%%%%%%%%%%%%%%%%%%%%%%%%%%%%%%%%%%%%%%%%%%%%%%%%%%%%%%%%%


% --| Math |-------------------------------------------------------

% -- Replacements --
\newcommand{\comp}{\ast}
%\renewcommand{\dagger}{+}


% -- new commands --
\providecommand{\abs}[1]{\lvert#1\rvert}
\providecommand{\Abs}[1]{\left\lvert#1\right\rvert}
\providecommand{\norm}[1]{\left\Vert#1\right\Vert}
\providecommand{\Trace}[1]{\ensuremath{\Tr\{\,#1\,\}}} % Trace /Spur
%

% -- differentials --
\renewcommand{\d}{\partial\mspace{2mu}} % partial diff
\newcommand{\td}{\,\mathrm{d}}	% total diff
\newcommand{\ddt}[1]{\frac{\td #1}{\td t}}

% -- Abbrevitations --
\renewcommand{\Re}{\text{Re}}			% Real value
\renewcommand{\Im}{\text{Im}}			% Real value
\newcommand{\complex}{\mathbb{C}} % Complex
\newcommand{\real}{\mathbb{R}}    % Real
%\newcommand{\R}{\real}						% Real
%\newcommand{\N}{\mathbb{N}}
%\newcommand{\Z}{\mathbb{Z}}
\renewcommand{\L}{\mathcal{L}}
\newcommand{\N}{\mathcal{N}}
\newcommand{\R}{\mathcal{R}}
\newcommand{\D}{\mathcal{D}}
%
\newcommand{\Ham}{\mathcal{H}}    % Hamilton
\newcommand{\Prob}{\mathscr{P}}    % Hamilton
\newcommand{\unity}{\mathds{1}}   % Real
%

\newcommand\gammab{\gamma_\bot}
\newcommand\gammap{\gamma_\parallel}
\newcommand\gammai{\gamma_\text{int}}
\newcommand\gammae{\gamma_\text{ext}}

% -- New Operators --
\DeclareMathOperator{\rot}{rot}
\DeclareMathOperator{\grad}{grad}
%\DeclareMathOperator{\div}{div}
\renewcommand{\div}{\text{div}\,}
\DeclareMathOperator{\Tr}{Tr}
\DeclareMathOperator{\const}{const}
\DeclareMathOperator{\e}{e} 			% exponatial Function

% -- new symbols --
\newcommand{\laplace}{\Delta}
\newcommand{\dalembert}{\Box}

% -- new arrows --
\renewcommand{\leadsto}{\Longrightarrow}
\newcommand{\leftrightleadsto}{\Longleftrightarrow}


% -- Text subscripts--
\newcommand{\rel}{_\text{rel}}
%\newcommand{\st}{\text{st}}
%

% -- other --
\newcommand{\com}[2]{\underbrace{#1}_{\textrm{\scriptsize #2}}}
\newcommand{\with}[1]{\stackrel{\ref{#1}}{\Longrightarrow}}
%\newcommand{\unit}[1]{\,\textrm{#1}}

%\newcommand{\variance}[1]{\delta \mean{#1}^2}
\newcommand{\variance}[1]{(\Delta{#1})^2}
%\newcommand{\variance}[1]{\delta #1^2}

% -- Physics --------------------------------
\newcommand\op[1]{{\hat{\mathrm{#1}}}}  % Operator

\newcommand\expect[1]{\ensuremath{\left\langle{#1}\right\rangle}} %
%
\newcommand{\mean}[1]{\ensuremath{\overline{#1}}} % mean value
%
\newcommand{\state}[1]{\ensuremath{\ket{#1}}}
%
\newcommand\commutator[2]{\ensuremath{\mathinner{%
    \mathopen[\,#1,#2\,\mathclose]}}}
\newcommand{\Commutator}[2]{\ensuremath{\left[\,#1,#2\,\right]}}
\newcommand{\bigcommutator}[2]{\ensuremath{\bigl[\,#1,#2\,\bigr]}}
\newcommand{\Bigcommutator}[2]{\ensuremath{\Bigl[\,#1,#2\,\Bigr]}}
%
\newcommand\poisson[2]{\mathinner{%
    \mathopen\{#1,#2\mathclose\}}}
%

% -- Layout --------------------------------

\newcommand*{\dashfill}{\leavevmode\cleaders\hbox{-}\hfill\kern0pt}

\newcommand*{\midhrulefill}{
\leavevmode
\cleaders\hbox to 1ex{\raisebox{.5ex}{\rule{1ex}{.4pt}}}\hfill\kern0pt
}


%% Kommandos fuer Tabellen. Entnommen aus The LateX Companion, tabsatz.ps und diversen Dokus:

%%% ---| Farben fuer Tabellen |-------------------
\IfPackageLoaded{xcolor}{
   \colorlet{tablesubheadcolor}{gray!30}
   \colorlet{tableheadcolor}{gray!25}
   \colorlet{tableblackheadcolor}{black!100}
   \colorlet{tablerowcolor}{gray!10.0}
}
%%% ---------------------------------------------


%%% -| Neue Spaltendefinitionen 'columntypes' |--
%
% Belegte Spaltentypen:
% l - links
% c - zentriert
% r - rechts
% p,m,b  - oben, mittig, unten
% X - tabularx Auto-Spalte

% um Tabellenspalten mit Flattersatz zu setzen, muss \\ vor
% (z.B.) \raggedright geschuetzt werden:
\newcommand{\PreserveBackslash}[1]{\let\temp=\\#1\let\\=\temp}


% Spalten mit Flattersatz und definierte Breite:
% m{} -> mittig
% p{} -> oben
% b{} -> unten
%
% Linksbuendig:
\newcolumntype{v}[1]{>{\PreserveBackslash\RaggedRight\hspace{0pt}}p{#1}}
\newcolumntype{M}[1]{>{\PreserveBackslash\RaggedRight\hspace{0pt}}m{#1}}
% % Rechtsbuendig :
% \newcolumntype{R}[1]{>{\PreserveBackslash\RaggedLeft\hspace{0pt}}m{#1}}
% \newcolumntype{S}[1]{>{\PreserveBackslash\RaggedLeft\hspace{0pt}}p{#1}}
% % Zentriert :
% \newcolumntype{Z}[1]{>{\PreserveBackslash\Centering\hspace{0pt}}m{#1}}
% \newcolumntype{A}[1]{>{\PreserveBackslash\Centering\hspace{0pt}}p{#1}}

\newcolumntype{Y}{>{\PreserveBackslash\RaggedLeft\hspace{0pt}}X}
%%% Spalten fuer Mathematik
%
% serifenlose Matheschrift
\newcolumntype{s}[1]{%
	>{\DC@{.}{,}{#1}\mathsf\bgroup}l%
	<{\egroup\DV@end}%
}

% Tabellenspaltentyp fuer den Kopf: (Farbe + Ausrichtung)
\newcolumntype{H}[1]{>{\columncolor{tableheadcolor}}l}

% aequivalent aus typokurz (fett+grau+links)
% \newcolumntype{H}{>{\fontseries{b}\selectfont%
%     \columncolor[gray]{.8}[6pt][0pt]}l}
%%% --------------------------------------------


%%% ---|Listen in Tabellen |--------------------
\newcommand{\removeindentation}{%
	\leftmargini=\labelsep%
	\advance\leftmargini by \labelsep%
}
%
\makeatletter
\newcommand\tableitemize{
	\@minipagetrue%
	\removeindentation
}
\makeatother
%%% --------------------------------------------

%%% ---|Layout der Tabellen |-------------------

% Neue Umgebung fuer Tabellen:

\newenvironment{Tabelle}[2][c]{%
  \tablestylecommon
  \begin{longtable}[#1]{#2}
  }
  {\end{longtable}%
  \tablerestoresettings
}


% Groesse der Schrift in Tabellen
\newcommand{\tablefontsize}{ \footnotesize}
\newcommand{\tableheadfontsize}{\footnotesize}

% Layout der Tabelle: Ausrichtung, Schrift, Zeilenabstand
\newcommand\tablestylecommon{%
  \renewcommand{\arraystretch}{1.4} % Groessere Abstaende zwischen Zeilen
  \normalfont\normalsize            %
  \sffamily\tablefontsize           % Serifenlose und kleine Schrift
  \centering%                       % Tabelle zentrieren
}

\newcommand{\tablestyle}{
	\tablestylecommon
	%\tablealtcolored
}

% Ruecksetzten der Aenderungen
\newcommand\tablerestoresettings{%
  \renewcommand{\arraystretch}{1}% Abstaende wieder zuruecksetzen
  \normalsize\rmfamily % Schrift wieder zuruecksetzen
}

% Tabellenkopf: Serifenlos+fett+schraeg+Schriftfarbe
\newcommand\tablehead{%
  \tableheadfontsize%
  \sffamily\bfseries%
  %\slshape
  %\color{white}
}

\newcommand\tablesubheadfont{%
  \tableheadfontsize%
  \sffamily\bfseries%
  \slshape
  %\color{white}
}


\newcommand\tableheadcolor{%
	%\rowcolor{tablesubheadcolor}
	%\rowcolor{tableblackheadcolor}
	\rowcolor{tableheadcolor}%
}

\newcommand\tablesubheadcolor{%
	\rowcolor{tablesubheadcolor}
	%\rowcolor{tableblackheadcolor}
}


\newcommand{\tableend}{\arrayrulecolor{black}\hline}

% Tabellenkopf (1=Spaltentyp, 2=Text)
% \newcommand{\tablehead}[2]{
%   \multicolumn{1}{#1@{}}{%
%     \raisebox{.1mm}{% Ausrichtung der Beschriftung
%       #2%
%     }\rule{0pt}{4mm}}% unsichtbare Linie, die die Kopfzeile hoeher macht
% }


\newcommand{\tablesubhead}[2]{%
  \multicolumn{#1}{>{\columncolor{tablesubheadcolor}}l}{\tablesubheadfont #2}%
}

% Tabellenbody (=Inhalt)
\newcommand\tablebody{%
\tablefontsize\sffamily\upshape%
}

\newcommand\tableheadshaded{%
	\rowcolor{tableheadcolor}%
}
\newcommand\tablealtcolored{%
	\rowcolors{1}{tablerowcolor}{white!100}%
}
%%% --------------------------------------------

\newlength{\mylen}
\newlength{\adjusthspace}

\newenvironment{tabularc}[2]
{%
	\setlength\mylen{#2/(#1)-\tabcolsep*2-\arrayrulewidth*(#1+1)/(#1)}%
	%\setlength{\adjusthspace}{((#2-1)/2)*\linewidth}
	%\par\noindent
	%\hspace*{-\the\adjusthspace}
	\begin{tabular}%{#2}%
		{*{#1}{v{\the\mylen}}}%
}
{\end{tabular}\par}


%%% Silbentrennung
\hyphenation{}

%% Dokument Beginn %%%%%%%%%%%%%%%%%%%%%%%%%%%%%%%%%%%%%%%%%%%%%%%%%%%%%%%%

% - Deckblatt,
% - Inhaltsverzeichnis,
% - Hauptteil gegliedert z.B. in
%   Einleitung, Grundlagen, Experimente, Ergebnisse, Zusammenfassung
% - Literaturverzeichnis,
% - Abbildungsverzeichnis (ggf.),
% - Tabellenverzeichnis (ggf.),
% - Abk�rzungsverzeichnis (ggf.),
% - Formelverzeichnis (ggf.),
% - Anhang, (nicht mehr Bestandteil der Arbeit! Wird daher nicht bewertet)
% - Erkl�rung der Urheberschaft,

\begin{document}
% Deckblatt
% \subject{Diplomarbeit \\ Universit�t <einf�gen>}
% \title{<Titel einf�gen>}
% \author{<Autor einf�gen>}
% \date{<Datum einf�gen>}
% \maketitle

% \begin{titlepage}
% 	\mbox{}\vspace{5\baselineskip}\\
% 	\sffamily\huge
% 	\centering
% 	<Titel einf�gen>
% 	\vspace{3\baselineskip}\\
% 	\rmfamily\Large
% 	Diplomarbeit \\ Universit�t <einf�gen>
% 	\vspace{2\baselineskip}\\
% 	\rmfamily\Large
% 	<Autor einf�gen>
% 	\vspace{1\baselineskip}\\
% 	<Datum einf�gen>
% \end{titlepage}


\begin{titlepage}
   \mbox{}\vspace{5\baselineskip}\\
   \sffamily\huge
   \centering
   % Titel
   \LaTeXe{}-Vorlage von Matthias Pospiech
   \vspace{3\baselineskip}\\
   \rmfamily\Large
   Leibniz Universit�t Hannover
   \vspace{2\baselineskip}\\
   \rmfamily\Large
   Matthias Pospiech
   \vspace{1\baselineskip}\\
   \today
\end{titlepage}


\thispagestyle{empty}
\chapter*{Erkl�rung der Selbstst�ndigkeit}
\thispagestyle{empty}
Hiermit versichere ich, die vorliegende Arbeit selbstst�ndig verfasst und keine anderen als die angegebenen Quellen und Hilfsmittel benutzt sowie die Zitate deutlich kenntlich gemacht zu haben.
\vspace{4\baselineskip}\\
<Ort einf�gen>, den <Datum einf�gen> \hfill <Autor einf�gen>
\vspace{4\baselineskip}\\
\clearpage
\mbox{}\thispagestyle{empty}

\frontmatter

\cleardoublepage
% Inhaltsverzeichnis in den PDF-Links eintragen
% \pdfbookmark[1]{Inhaltsverzeichnis}{toc}
\tableofcontents

% Hauptteil
\mainmatter
% Demonstration der Vorlage
% \newcommand{\env}[1]{\texttt{#1}}
\newcommand{\command}[1]{\texttt{#1}}
\newcommand{\package}[1]{\texttt{\itshape#1}}
\newcommand{\engl}[1]{(engl: \textit{#1})\xspace}

\IfDefined{chapter}{%
\chapter{Demonstration eines Dokumentes}
}

\IfDefined{minitoc}{%
\minitoc
}

% ============================================================
\section{Textauszeichnungen}

Die Standard Text Schalter erzeugen
\textbf{fett},
\textit{kursiv\footnote{kursiv und schr�ggestellt sind unterschiedliche Schriften und daher nicht das gleiche !} (italic)},
\textsl{schr�ggestellt (slanted)},
\textsf{serifenlos (grotesk)},
\textsc{Kapit�lchen} und
\texttt{Schreibmaschinenschrift}.
Sowie beliebige Kombinationen derselben:
\textit{\textbf{fett kursiv}},
\textsl{\textbf{fett schr�g}},
\textsf{\textbf{fett serifenlos}},
\textsc{\textbf{Fette Kapit�lchen}}
und
\textsl{\textsf{serifenlos schr�g}}. Je nach Schrift sind jedoch nicht alle Kombinationen m�glich. In dem Fall bekommt man die Fehlermeldung `Some font shapes were not available, defaults substituted.'

\IfDefined{so}{%
Mit dem Paket \package{soul} erh�lt man zus�tzlich
\so{Sperrung},
\ul{Unterstreichen} mit der M�glichkeit des Zeilenumbruchs,
\st{durchstreichen}, sowie
\hl{hervorheben}. Letzteres wird allerdings in dvi nicht korrekt angezeigt.
}
\IfDefined{uuline}{%
Mit dem Paket \package{ulem} erh�lt man noch weitere M�glichkeiten des Unterstreichens:
\uline{einfach untersrichen},
\uuline{doppelt untersrichen},
\uwave{wellig untermalt},
\sout{durchgestrichen} und
\xout{deutlicher durch�gestrichen}.
}
\IfDefined{url}{%
Etwas spezieller wird es beim Erzeugen von Webseiten Links (URLs). Das Hauptproblem hierbei ist das Umbrechen solcher. Hier ein Beispiel: \url{http://www.dante.de}. Die daf�r genutzte Schrift kann in der Pr�ambel eingestellt werden.
}
% Darstellen von Code, insbesondere \LaTeX-Code z.\,B. |\textbf{fett}| innerhalb von \LaTeX~mit dem Paket \package{fancyvrb}

\IfDefined{enquote}{%
Zitieren, auch doppelt geht sehr einfach mit dem Paket \package{csquote}. Hier ein Unsinniger Satz: \enquote{Eine Aussage die eine zweite Aussage: \enquote{mit ge�ndertem Anf�hrungszeichen} darstellt}.
}
% ============================================================

\section{Text und Zitate}
% ------------------------------------------------------------
\IfDefined{FlushLeft}{%
\subsection{Linksb�ndig}
%
Mit dem Paket \package{ragged2e} und der Umgebung \env{FlushLeft} bzw. dem Befehl \command{Raggedright} wird Flattersatz (Gegenteil von Blocksatz) gegen�ber dem \LaTeX{} Standard noch verbessert.

\begin{FlushLeft}
Aliquam ultrices libero hendrerit diam. Vestibulum ultrices sapien sit amet elit. Quisque tempor nisl eu sem. Nam lorem lectus, viverra nec, rutrum quis, lobortis nec, magna. Praesent hendrerit tortor vitae elit. Vivamus sed leo at mi elementum semper. Lorem ipsum dolor sit amet, consectetuer adipiscing elit. Aliquam eu nisi. Nam eget dui a tortor congue imperdiet. Etiam mattis. Nam tristique. Sed malesuada neque ut leo. Aenean est. In id augue.
\end{FlushLeft}
} %Endif
% ------------------------------------------------------------
\IfDefined{FlushRight}{%
\subsection{Rechtsb�ndig}
%
Mit dem Paket \package{ragged2e} und der Umgebung \env{FlushRight} bzw. dem Befehl \command{RaggedLeft} wird Flattersatz (Gegenteil von Blocksatz) gegen�ber dem \LaTeX{} Standard noch verbessert.

\begin{FlushRight}
Aliquam ultrices libero hendrerit diam. Vestibulum ultrices sapien sit amet elit. Quisque tempor nisl eu sem. Nam lorem lectus, viverra nec, rutrum quis, lobortis nec, magna. Praesent hendrerit tortor vitae elit. Vivamus sed leo at mi elementum semper. Lorem ipsum dolor sit amet, consectetuer adipiscing elit. Aliquam eu nisi. Nam eget dui a tortor congue imperdiet. Etiam mattis. Nam tristique. Sed malesuada neque ut leo. Aenean est. In id augue.
\end{FlushRight}
%
} % Endif
% ------------------------------------------------------------
\IfDefined{multicols}{%
\subsection{Mehrere Spalten}
%
\begin{multicols}{3}[Text mit 3 Spalten. Erstellt mit dem Paket \package{multicol}]
Suspendisse ac nibh vitae nunc iaculis accumsan. Vivamus venenatis, orci vitae interdum tristique, nisl lectus fermentum arcu, sed vehicula pede orci et nunc. Cras tempus ultrices leo. Nulla at tortor. Morbi nisl tellus, lobortis nec, nonummy a, vulputate at, felis. In interdum varius sem. Fusce pellentesque, eros vitae consectetuer dignissim, ipsum urna tincidunt urna, ut aliquet libero lectus vel purus. In commodo iaculis justo. Sed euismod. Praesent molestie leo ac erat. Etiam a felis.
% Nunc ipsum diam, porta ac, mollis non, mattis a, felis. Etiam nisl sapien, malesuada eget, rutrum at, dictum non, metus. Aliquam ut nunc in purus rutrum posuere. Proin id risus. Integer dignissim, lorem sit amet cursus adipiscing, sapien purus posuere erat, ac porta risus augue non enim. Fusce nunc nunc, sodales et, vestibulum ut, auctor ac, sem. Vivamus nisi lectus, consectetuer eget, congue at, feugiat et, elit. Praesent sem. Curabitur interdum placerat odio.
\end{multicols}
} % Endif
% ------------------------------------------------------------
\subsection{Zentierter Text}
%
Zentierter Text, erstellt mit \env{center}\footnote{nicht verwenden innerhalb \env{table} oder \env{figure}!}.
\begin{center}
Aliquam ultrices libero hendrerit diam. Vestibulum ultrices sapien sit amet elit. Quisque tempor nisl eu sem. Nam lorem lectus, viverra nec, rutrum quis, lobortis nec, magna. Praesent hendrerit tortor vitae elit. Vivamus sed leo at mi elementum semper. Lorem ipsum dolor sit amet, consectetuer adipiscing elit. Aliquam eu nisi. Nam eget dui a tortor congue imperdiet. Etiam mattis. Nam tristique. Sed malesuada neque ut leo. Aenean est. In id augue.
\end{center}

% ------------------------------------------------------------
\subsection{Zitat mit quote}
Ein Text in einer \env{quote} Umgebung\footnote{F�r Zitate eignet sich allerdings eher die Umgebung \env{blockquote} vom Paket \package{csquotes}}. Der gesammte Text innerhalb der Umgebung wird von beiden Seiten einger�ckt.

\begin{quote}
Aliquam ultrices libero hendrerit diam. Vestibulum ultrices sapien sit amet elit. Quisque tempor nisl eu sem. Nam lorem lectus, viverra nec, rutrum quis, lobortis nec, magna. Praesent hendrerit tortor vitae elit. Vivamus sed leo at mi elementum semper. Lorem ipsum dolor sit amet, consectetuer adipiscing elit. Aliquam eu nisi. Nam eget dui a tortor congue imperdiet. Etiam mattis. Nam tristique. Sed malesuada neque ut leo. Aenean est. In id augue.
\end{quote}

% ------------------------------------------------------------
\IfDefined{blockquote}{%
%
\subsection{Zitat mit blockquote}
Nam dui ligula, fringilla a, euismod sodales,
sollicitudin vel, wisi. Morbi auctor lorem non justo. Nam lacus
libero, pretium at, lobortis vitae, ultricies et, tellus. Donec
aliquet, tortor sed accumsan bibendum, erat ligula aliquet magna,
vitae ornare odio metus a mi.
% Morbi ac orci et nisl hendrerit mollis. Suspendisse ut massa. Cras nec ante. Pellentesque a nulla. Cum sociis natoque penatibus et magnis dis parturient montes, nascetur ridiculus mus. Aliquam tincidunt urna. Nulla ullamcorper vestibulum turpis. Pellentesque cursus luctus mauris.

\blockquote[(Lorem Lipsum, P. 50)]{Lorem ipsum dolor sit amet, consectetur adipisicing elit, sed do eiusmod tempor incididunt ut labore et dolore magna aliqua. Ut enim ad minim veniam, quis nostrud exercitation ullamco laboris nisi ut aliquip ex ea commodo consequat. Duis aute irure dolor in reprehenderit in voluptate velit esse cillum dolore eu fugiat nulla pariatur.}
} %Endif
% ------------------------------------------------------------
\subsection{Zitat mit mit \command{cite}}
Je nach verwendetem Paket zum Zitieren landen die Zitate direkt im Text (\package{natbib}) oder in den Fussnote bei \package{jurabib} mit der Option `super'.

\minisec{Aufbau von Glasfasern}
Eine optische Faser besteht in der Mitte aus dem Kern \engl{core} welcher von einem Mantel \engl{cladding} umgeben ist. Der Licht f�hrende Teil ist der Kern. Entscheidend f�r die Lichtwellenleitung ist der Brechungsindexsprung vom Licht f�hrenden Kern zum Mantel.
Wegen ihres Aufbaus werden diese Fasern auch \emph{Stufenindexfasern}\footnote{Es existieren noch weiter Typen  von Glasfasern, f�r Faserlaser spielen diese jedoch nur eine geringe Bedeutung und werden daher auch nicht weiter betrachtet. Eine Klassifikation findet man in \cite{Yeh}} genannt.
Sowohl der Kern als auch der Mantel werden typischerweise aus Quarzglas gefertigt, wobei der h�here Brechungsindex im Kern durch zus�tzliche Dotierung mit GeO\textsubscript{2} oder P\textsubscript{2}O\textsubscript{5} erreicht wird \cite{Agrawal}.
Typische Brechungsindexzahlen f�r Kern und Mantel sind 1,47 und 1,46 \cite{GuidetofiberOptics}.
Umgeben wird eine Faser zus�tzlich von einer Beschichtung \engl{coating}, die als Schutz dient.
Die Beschichtung besteht aus einem Polymer, das die Faser vor �u�eren chemischen Einfl�ssen sch�tzt und die mechanische Stabilit�t verbessert.


% ============================================================
\section{Listen}

% ------------------------------------------------------------
\subsection{itemize}
Dies ist die Standard Aufz�hlungsliste von \LaTeXe. Sie hat einen Abstand zwischen den Eintr�gen um bei L�ngeren Zeilen das Lesen zu erleichtern.

\begin{itemize}
   \item Lorem ipsum dolor sit amet, consectetur adipisicing elit, sed do eiusmod tempor incididunt ut labore et dolore magna aliqua.
%
   \item Lorem ipsum dolor sit amet, consectetur adipisicing elit, sed do eiusmod tempor incididunt ut labore et dolore magna aliqua.
%
   \item Lorem ipsum dolor sit amet, consectetur adipisicing elit, sed do eiusmod tempor incididunt ut labore et dolore magna aliqua.
\end{itemize}

mit mehr als einer Ebene

\begin{itemize}
   \item Lorem ipsum dolor sit amet, consectetur adipisicing elit, sed do eiusmod tempor incididunt ut labore et dolore magna aliqua.
%
   \begin{itemize}
      \item Lorem ipsum dolor sit amet, consectetur adipisicing elit, sed do eiusmod tempor incididunt ut labore et dolore magna aliqua.
      %
      \begin{itemize}
         \item Lorem ipsum dolor sit amet, consectetur adipisicing elit, sed do eiusmod tempor incididunt ut labore et dolore magna aliqua.
         %
         \item Lorem ipsum dolor sit amet, consectetur adipisicing elit, sed do eiusmod tempor incididunt ut labore et dolore magna aliqua.
         %
         \item Lorem ipsum dolor sit amet, consectetur adipisicing elit, sed do eiusmod tempor incididunt ut labore et dolore magna aliqua.
      \end{itemize}
      %
      \item Lorem ipsum dolor sit amet, consectetur adipisicing elit, sed do eiusmod tempor incididunt ut labore et dolore magna aliqua.
      %
      \item Lorem ipsum dolor sit amet, consectetur adipisicing elit, sed do eiusmod tempor incididunt ut labore et dolore magna aliqua.
   \end{itemize}
\end{itemize}

% ------------------------------------------------------------
\subsection{enumerate}
Dies ist die Standard Nummerierungsliste von \LaTeXe{}. Sie hat einen Abstand zwischen den Eintr�gen um bei L�ngeren Zeilen das Lesen zu erleichtern.

\begin{enumerate}
   \item Lorem ipsum dolor sit amet, consectetur adipisicing elit, sed do eiusmod tempor incididunt ut labore et dolore magna aliqua.
%
   \begin{enumerate}
      \item Lorem ipsum dolor sit amet, consectetur adipisicing elit, sed do eiusmod tempor incididunt ut labore et dolore magna aliqua.
      %
      \begin{enumerate}
         \item Lorem ipsum dolor sit amet, consectetur adipisicing elit, sed do eiusmod tempor incididunt ut labore et dolore magna aliqua.
         %
         \item Lorem ipsum dolor sit amet, consectetur adipisicing elit, sed do eiusmod tempor incididunt ut labore et dolore magna aliqua.
         %
         \item Lorem ipsum dolor sit amet, consectetur adipisicing elit, sed do eiusmod tempor incididunt ut labore et dolore magna aliqua.
      \end{enumerate}
      %
      \item Lorem ipsum dolor sit amet, consectetur adipisicing elit, sed do eiusmod tempor incididunt ut labore et dolore magna aliqua.
      %
      \item Lorem ipsum dolor sit amet, consectetur adipisicing elit, sed do eiusmod tempor incididunt ut labore et dolore magna aliqua.
   \end{enumerate}
\end{enumerate}

% ------------------------------------------------------------
%
\subsection{Kompakte Liste}

Die eleganteste L�sung bietet das Paket \package{enumitem} mit der Option `noitemsep'
\begin{itemize}[noitemsep]
\item Diese Umgebung
\item sollte man nur nutzen,
\item wenn die Eintr�ge nicht l�nger als eine Zeile sind
\end{itemize}

% ------------------------------------------------------------
\subsection{Beliebige Labels}

Dies ist die \env{enumerate} Umgebung des Paketes \package{enumitem} mit der Option [label=\alph{enumi})].

\begin{enumerate}[label=\alph{enumi})]
         \item Lorem ipsum dolor sit amet, consectetur adipisicing elit, sed do eiusmod tempor incididunt ut labore et dolore magna aliqua.
         %
         \item Lorem ipsum dolor sit amet, consectetur adipisicing elit, sed do eiusmod tempor incididunt ut labore et dolore magna aliqua.
\end{enumerate}

% ------------------------------------------------------------
\subsection{description}
%
\begin{description}
\item[Beschreibung] Die Umgebung \env{description} dient zum Erl�utern.
\item[Fl�sse] Elbe, Weser
\item[Meere] Indischer, Pazifischer Ozean, Mittelmeer
\end{description}

% ============================================================
\section{Abbildungen}

% ------------------------------------------------------------
\IfDefined{includegraphics}{%
%
\subsection{Normale figure Umgebung}
Nulla malesuada porttitor diam. Donec felis
erat, congue non, volutpat at, tincidunt tristique, libero. Vivamus
viverra fermentum felis. Donec nonummy pellentesque ante. Phasellus
adipiscing semper elit. Proin fermentum massa ac quam. Sed diam
turpis, molestie vitae, placerat a, molestie nec, leo. Maecenas
lacinia.
%
\begin{figure}[H]
   \centering
   \includegraphics[width=0.3\textwidth]{images/latex}
   \caption[Normale Figure Umgebung]{Dies ist eine lange Abbildungungsbeschriftung. Dies ist eine lange Abbildungungsbeschriftung. Dies ist eine lange Abbildungungsbeschriftung.}
\end{figure}

\subsection{Normale figure Umgebung, zwei Bilder nebeneinander}
\begin{figure}[H]
   \centering
   \begin{minipage}{0.4\linewidth}
   \includegraphics[width=1.0\linewidth]{images/latex}
   \caption[Normale Figure Umgebung]{Dies ist eine lange Abbildungungsbeschriftung. Dies ist eine lange Abbildungungsbeschriftung. Dies ist eine lange Abbildungungsbeschriftung.}
	\end{minipage}\hspace{2em}%\hfill%
   \begin{minipage}{0.4\linewidth}
   \includegraphics[width=1.0\linewidth]{images/latex}
   \caption[Normale Figure Umgebung]{Dies ist eine lange Abbildungungsbeschriftung. Dies ist eine lange Abbildungungsbeschriftung. Dies ist eine lange Abbildungungsbeschriftung.}
	\end{minipage}
\end{figure}

%
} % Endif
% ------------------------------------------------------------
\IfDefined{subfloat}{%
%
\subsection{Normale figure Umgebung mit subfloat}
Nam ipsum ligula, eleifend at, accumsan nec, suscipit a,
ipsum. Morbi blandit ligula feugiat magna. Nunc eleifend consequat
lorem. Sed lacinia nulla vitae enim. Pellentesque tincidunt purus
vel magna. Integer non enim. Praesent euismod nunc eu purus. Donec
bibendum quam in tellus. Nullam cursus pulvinar lectus. Donec et mi.
Nam vulputate metus eu enim. Vestibulum pellentesque felis eu
massa.
%
\begin{figure}[H]
   \centering
   \subfloat[Erste Bild]{\includegraphics[width=0.3\textwidth]{images/latex}}
      \hspace{0.1\textwidth}
   \subfloat[Zweites Bild]{\includegraphics[width=0.3\textwidth]{images/latex}}
   \caption{Ein Beispiel f�r subfigures}
\end{figure}

} % Endif
% ------------------------------------------------------------
\IfDefined{captionof}{%
%
\subsection{captionof}

Pellentesque et lectus a est imperdiet egestas. Fusce tempus facilisis lacus. Morbi porttitor eleifend dolor. Integer ante. Integer ornare. Quisque ac urna. Nam egestas, eros sed tempor tincidunt, dolor ante dapibus tellus, eget dapibus metus augue et lectus.

\begin{center}
   \includegraphics[width=0.3\textwidth]{images/latex}
   \captionof{figure}{Ein Beispiel f�r ein nichtflie�endes Bild mit Caption durch `captionof'}
\end{center}

Nullam sed enim quis enim blandit commodo. Nulla ultricies metus eu lorem. Phasellus tincidunt ullamcorper orci. Cras lectus metus, luctus id, vulputate quis, tempor non, nulla. Nullam ut eros et elit faucibus viverra. Cras eu turpis eu urna ullamcorper feugiat. Aliquam bibendum rhoncus tortor. Etiam eu orci. In congue congue quam. Maecenas nec sem quis mi ultrices luctus. Aliquam et risus eu eros molestie pellentesque. Aenean quis purus sit amet lorem porttitor consectetuer.
%
} %Endif

% ------------------------------------------------------------
\IfDefined{wrapfigure}{
%
\subsection{wrapfigure}
\begin{wrapfigure}{r}{0.3\textwidth}
	\includegraphics[width=0.3\textwidth]{images/latex}
	\caption{Ein Beispiel f�r wrapfigure }
   %\vspace{-2\baselineskip}
\end{wrapfigure}
Suspendisse vel felis. Ut lorem lorem, interdum
eu, tincidunt sit amet, laoreet vitae, arcu. Aenean faucibus pede eu
ante. Praesent enim elit, rutrum at, molestie non, nonummy vel,
nisl. Ut lectus eros, malesuada sit amet, fermentum eu, sodales
cursus, magna. Donec eu purus. Quisque vehicula, urna sed ultricies
auctor, pede lorem egestas dui, et convallis elit erat sed nulla.
Donec luctus. Curabitur et nunc. Aliquam dolor odio, commodo
pretium, ultricies non, pharetra in, velit. Integer arcu est,
nonummy in, fermentum faucibus, egestas vel, odio.
} %Endif
% ------------------------------------------------------------
\IfDefined{wrapfloat}{
%
\subsection{wrapfloat}
%
\begin{wrapfloat}{figure}{r}{0.3\textwidth}
	\includegraphics[width=0.3\textwidth]{images/latex}
	\caption{Ein Beispiel f�r wrapfloat }
\end{wrapfloat}
Ut vulputate lorem a ante. Praesent risus purus, feugiat ac, dignissim at, scelerisque nec, tortor. Nulla a ipsum. Integer ut lorem. Cras convallis pulvinar velit. Pellentesque habitant morbi tristique senectus et netus et malesuada fames ac turpis egestas. Vestibulum tempus. Fusce suscipit semper est. Vestibulum ante ipsum primis in faucibus orci luctus et ultrices posuere cubilia Curae; Aliquam sagittis pretium dui. Cum sociis natoque penatibus et magnis dis parturient montes, nascetur ridiculus mus. In velit. Nunc porta feugiat ipsum. Sed dolor. Phasellus porttitor mattis libero. Nam justo mi, tincidunt vitae, mattis sit amet, dapibus et, mi. Nunc in nisi eu lacus ultricies feugiat. Donec volutpat posuere arcu.

} % Endif

% ------------------------------------------------------------
%
\IfDefined{floatingfigure}{

\subsection{floatingfigure}
%
\begin{floatingfigure}[r]{0.3\textwidth}
   \includegraphics[width=0.3\textwidth]{images/latex}
   \caption{Ein Beispiel f�r floatingfigure}
\end{floatingfigure}

Ut vulputate lorem a ante. Praesent risus purus, feugiat ac, dignissim at, scelerisque nec, tortor. Nulla a ipsum. Integer ut lorem. Cras convallis pulvinar velit. Pellentesque habitant morbi tristique senectus et netus et malesuada fames ac turpis egestas. Vestibulum tempus. Fusce suscipit semper est. Vestibulum ante ipsum primis in faucibus orci luctus et ultrices posuere cubilia Curae; Aliquam sagittis pretium dui. Cum sociis natoque penatibus et magnis dis parturient montes, nascetur ridiculus mus. In velit. Nunc porta feugiat ipsum. Sed dolor. Phasellus porttitor mattis libero. Nam justo mi, tincidunt vitae, mattis sit amet, dapibus et, mi. Nunc in nisi eu lacus ultricies feugiat. Donec volutpat posuere arcu.



} % Endif
% ------------------------------------------------------------
\IfDefined{wrapfloat}{

%
\subsection{wrapfloat mit subfloat}

\begin{wrapfloat}{figure}{r}{0.22\textwidth}
   \centering
   \vspace{-\baselineskip}
   \subfloat[erste Bild]%
      {\includegraphics[width=0.22\textwidth]{images/latex}}
   \vspace{1em}
   \subfloat[zweite Bild]%
      {\includegraphics[width=0.22\textwidth]{images/latex}}
   \caption{zwei Bilder}
    \vspace{-1\baselineskip}
\end{wrapfloat}

Pellentesque mollis nunc sed mauris tempor molestie. Aliquam adipiscing nisi eu metus. Proin viverra odio ac lorem consequat condimentum. Suspendisse bibendum tellus. Duis non diam. Aliquam sodales sapien in mauris. Sed euismod adipiscing lorem. Pellentesque nulla augue, nonummy vel, tincidunt at, blandit non, erat. Ut sodales, est non semper placerat, velit est lacinia orci, quis auctor nisl elit vel mi. Vivamus elementum turpis in ligula. Nunc dui sem, mattis et, rutrum sit amet, iaculis eu, urna. Praesent quis nunc. Class aptent taciti sociosqu ad litora torquent per conubia nostra, per inceptos hymenaeos. Suspendisse justo purus, lacinia et, iaculis ac, sagittis at, enim. Nulla placerat, massa iaculis imperdiet tristique, dolor libero laoreet risus, ac sollicitudin enim sem vel ante. Nam dignissim sem eu dui. Fusce et leo eu est vehicula ultricies. Aenean laoreet, enim vel bibendum malesuada, mauris est dictum velit, vitae suscipit augue ligula pharetra enim.

} % Endif
% ------------------------------------------------------------
 \IfDefined{SCfigure}{
\subsection{SCfigure}

\begin{SCfigure}[5][H]
 \includegraphics[width=0.4\textwidth]{images/latex}
 \caption[Side Caption Figure]{Suspendisse justo purus, lacinia et, iaculis ac, sagittis at, enim. Nulla placerat, massa iaculis imperdiet tristique, dolor libero laoreet risus, ac sollicitudin enim sem vel ante. Nam dignissim sem eu dui. Fusce et leo eu est vehicula ultricies. Aenean laoreet, enim vel bibendum malesuada, mauris est dictum velit, vitae suscipit augue ligula pharetra enim.}
\end{SCfigure}
%
} % Endif
% ============================================================
\section{Tabellen}

Bei allen hier vorgestellten Tabellen sind Befehle in die Tabellen eingef�gt die den Stil ver�ndern.

%--Einstellungen f�r Tabellen ----------
\renewcommand\tablestylecommon{%
  \renewcommand{\arraystretch}{1.4} % Groessere Abstaende zwischen Zeilen
  \normalfont\normalsize            %
  \sffamily\tablefontsize           % Serifenlose und kleine Schrift
  \centering%                       % Tabelle zentrieren
}

\renewcommand{\tablestyle}{%
   \tablestylecommon%
}

\renewcommand\tablebody{%
   \tablefontsize\sffamily\upshape%
}
% ------------------------------------------------------------
\IfDefined{bottomrule}{%
%
\subsection{Tabelle im \emph{booktabs} Stil}
%
\begin{table}[H]
\tablestyle
\begin{tabular}{lll}
\toprule
   \tablehead Tabellenkopf &
   \tablehead Tabellenkopf &
   \tablehead Tabellenkopf \tabularnewline
\midrule
Inhalt & Inhalt  & Inhalt \tabularnewline
Inhalt & Inhalt  & Inhalt \tabularnewline
Inhalt & Inhalt  & Inhalt \tabularnewline
\bottomrule
\end{tabular}
\end{table}
%
} % Endif

% ------------------------------------------------------------
\subsection{Tabelle ohne Kopf}

\begin{table}[H]
\tablestyle
\begin{tabular}{lll}
\hline
\tablebody
Inhalt & Inhalt  & Inhalt \tabularnewline
Inhalt & Inhalt  & Inhalt \tabularnewline
Inhalt & Inhalt  & Inhalt \tabularnewline
\hline
\end{tabular}
\end{table}

\IfDefined{tabularx}{%
%
\subsection{Tabelle mit tabularx}
Diese Tabelle enth�lt zus�tzlichen Code f�r den Stil des Kopfes der Tabelle.
%--------------------------------------------------------
\begin{table}[H]
\tablestyle
\begin{tabularx}{\textwidth}{lXXlX}
\tableheadcolor
   \tablehead Tabellenkopf &
   \tablehead Tabellenkopf &
   \tablehead Tabellenkopf &
   \tablehead Tabellenkopf &
   \tablehead Tabellenkopf \tabularnewline
%
\tablebody
   \textit{Beschreibung} & Inhalt & Inhalt & Inhalt & Inhalt \tabularnewline
   \textit{Beschreibung} & Inhalt & Inhalt & Inhalt & Inhalt \tabularnewline
\tableend
\end{tabularx}
\caption{Tabelle mit tabularx}
\end{table}
%
} % Endif
% ------------------------------------------------------------
\IfDefined{rowcolor}{%
%
\subsection{Tabelle mit alternatierender Farbe}


%--Einstellungen f�r Tabellen ----------

\colorlet{tablesubheadcolor}{gray!40}
\colorlet{tableheadcolor}{gray!25}
\colorlet{tableblackheadcolor}{black!60}
\colorlet{tablerowcolor}{gray!15.0}


\renewcommand\tablehead{%
  \tableheadfontsize%
  \sffamily\bfseries%
  \slshape
  \color{white}
}

\renewcommand\tableheadcolor{
   \rowcolor{tableblackheadcolor}
}


%---------------------------------------
%
\begin{table}[H]
   \tablestyle
   \tablealtcolored
   \begin{tabular}{*{2}{v{0.45\textwidth}}}
   \hline
   \tableheadcolor
\tablehead Tabellenkopf &
\tablehead Tabellenkopf \tabularnewline\hline
% Zwischenkopf
\multicolumn{2}{>{\columncolor{tablesubheadcolor}}l}{
   \bfseries Zwischenkopf
} \tabularnewline

\tablebody
 Inhalt  & Inhalt \tabularnewline
 Inhalt  & Inhalt \tabularnewline
 Inhalt  & Inhalt \tabularnewline
 Inhalt  & Inhalt \tabularnewline
 Inhalt  & Inhalt \tabularnewline
 Inhalt  & Inhalt \tabularnewline
\multicolumn{2}{>{\columncolor{tablesubheadcolor}}l}{
   \bfseries Zwischenkopf
} \tabularnewline
 Inhalt  & Inhalt \tabularnewline
 Inhalt  & Inhalt \tabularnewline
   \hline
   \end{tabular}
\end{table}
%
} % End If
% ------------------------------------------------------------
\IfDefined{LTXtable}{%
\subsection{Seiten�bergreifende Tabellen}



%--Einstellungen f�r Tabellen ----------

\colorlet{tablerowcolor}{gray!10.0}

\renewcommand\tableheadcolor{
   \rowcolor{tableheadcolor}
}

\renewcommand\tablehead{%
  \tableheadfontsize%
  \sffamily\bfseries%
  \slshape
  \color{black}
}

%---------------------------------------

{
   \tablestyle
   \tablealtcolored
   \LTXtable{\textwidth}{tabellen/LongtableBeispiel.tex}
}

} % End If


% ============================================================

% \section{Quellcode (Verbatim Text)}
% ------------------------------------------------------------
%
% \subsection{fancyvrb}
% %
% Das Paket \package{fancyvrb} bietet Umfangreiche M�glichkeiten Code  darzustellen. Hier ein Beispiel:
% %
% \begin{Verbatim}[frame=single,framesep=5mm,numbers=left,rulecolor=\color{red}]
% \documentclass[a4paper,11pt,german]{scrreprt}
% \begin{document}
% Hello World
% \end{document}
% \end{Verbatim}
% %
% ------------------------------------------------------------
%
% \subsection{listings}
% %
% F�r Programm Listings nutzt man am besten das Paket \package{listings}. Hier ein einfaches Beispiel mit Visual Basic Code:
% \begin{lstlisting}[firstnumber=last]
% ' *************************************************
% ' Diagramme_entfernen()
% '
% ' Alle Diagramme auf dem aktuelles Blatt entfernen
% ' *************************************************
%
% Public Sub Diagramme_entfernen()
%
%   Dim Diagramm As Variant
%   Dim DiaNamen()
%   i = 0
%   'Nur wenn auch Diagramme vorhanden sind
%   If ActiveSheet.ChartObjects.Count() > 0 Then
%     ReDim Preserve DiaNamen(1 To ActiveSheet.ChartObjects.Count())
%     For Each Diagramm In ActiveSheet.ChartObjects
%       i = i + 1
%       DiaNamen(i) = Diagramm.name
%
%       ActiveSheet.ChartObjects(DiaNamen(i)).Activate
%       ActiveChart.ChartArea.Select
%       ActiveWindow.Visible = False
%       Selection.Delete
%     Next
%   End If
% End Sub
% \end{lstlisting}
% %
% ============================================================
% \section{Wissenschaftliches}

% ============================================================
\section{Mathematik}

\subsection{amsmath Umgebungen}

Im folgenden sind einige Beispiel f�r amsmath Umgebungen aufgelistet. Diese sind der Dokumentation `mathmode.pdf' von Herbert Voss entnommen.

\subsubsection{align}
\begin{align}
y      & = d\\
y      & = cx + d\\
y_{12} & = bx^{2} + cx + d\\
y(x)   & = ax^{3} + bx^{2}+ cx + d
\end{align}

\subsubsection{alignat}
\begin{alignat}{3}
abc   &= xxx &&= xxxxxxxxxxxx &&= aaaaaaaaa \\
ab    &= yyyyyyyyyyyyyyy &&= yyyy &&= ab
\end{alignat}

\subsubsection{multline}

\begin{multline}
a+b+c+d+e+f\\
+i+j+k+l+m+n
\end{multline}

\subsubsection{cases}
\begin{equation}
x = \begin{cases}
   0 & \text{if A =...}\\
   1 & \text{if B =...}\\
   x & \parbox{5 cm}{%
\flushleft %
this runs with as much text as you like,
but without an automatic linebreak,
it runs out of page ....} %
\end{cases}
\end{equation}

% ------------------------------------------------------------
\IfDefined{Braket}{%
\subsection{braket}
\begin{equation}
   \text{bra:} \Bra{a} \qquad \text{ket:} \Ket{a} \qquad \text{braket:} \Braket{a|b} \quad,\qquad \Braket{a|A|b}
\end{equation}
} % Endif
% ------------------------------------------------------------
\IfDefined{cancel}{%
\subsection{cancel}
\begin{equation}
f(x) = \frac{\cancel{(a+1)}x}{(x-1)\cancel{(a+1)}}
\end{equation}
} % Endif
% ------------------------------------------------------------
\IfDefined{empheq}{%
\subsection{empheq}
\begin{empheq}[box=\fbox]{align}
f(x) & = e^{-E/kT}
\end{empheq}
} % Endif

% ============================================================
\section{Index, Glossar, Symbolverzeichnis \& Co.}

% ------------------------------------------------------------
\IfDefined{index}{
%
\subsection{Text enh�lt Index}
Lorem\index{Lorem} ipsum\index{ipsum} dolor\index{dolor} sit amet, consectetuer adipiscing elit
Nam dui ligula, fringilla a, euismod sodales, sollicitudin vel, wisi.

} % Endif
% ------------------------------------------------------------
\IfDefined{nomenclature}{
%
\subsection{Symbole}
Es folgen einige Symbole die im Symbolverzeichnis auftauchen sollen: \\
\%  \&  �
\nomenclature{\%}{Prozentzeichnen}
\nomenclature{\&}{Kaufm�nisches Und}
\nomenclature{�}{Paragraph}

} % Endif
% ------------------------------------------------------------
\IfDefined{acs}{
%
\subsection{Abk�rzungsverzeichnis}

\minisec{Beispieltext aus der acronym Dokumentation}
In the early nineties, \acs{GSM} was deployed in many European
countries. \ac{GSM} offered for the first time international
roaming for mobile subscribers. The \acs{GSM}'s use of \ac{TDMA} as
its communication standard was debated at length. And every now
and then there are big discussion whether \ac{CDMA} should have
been chosen over \ac{TDMA}.

\begin{acronym}[TDMA]
 \acro{CDMA}{Code Division Multiple Access}
 \acro{GSM}{Global System for Mobile communication}
 \acro{TDMA}{Time Division Multiple Access}
\end{acronym}
} % Endif
% ------------------------------------------------------------


\chapter{Introduction}
\section{Personal motivation}
This thesis describes the analysis, enhanced design and implementation of an existing microcontroller based mesh solution \cite{korniowski}. The current solution showed.

\section{Research overview}

\chapter{Evaluation}
\section{Existing solution}
\section{Assumptions}
\section{Requirements}

\chapter{Hardware Design}
\section{RAM}
\begin{itemize}
\item Harvard architecture
\item RAM bus
\item Latch
\end{itemize}
\section{USB Serial Device}
\section{RFM12B Radio}
\section{Keyboard}

\chapter{Software Modules}
\section{UART}
\section{SPI}
\section{Watchdog}
\section{Timer}
\section{Shell}
\section{Network Stack}
\section{RFM12 Driver}

\chapter{Software Algorithms}
\section{Protothreads}
Designing a software system that executes on embedded micro-controllers implies a lot of challenges when many software modules are involved and complexity grows. The conceptually defined modules must be somehow implemented. If the micro-controller lacks an operating system then there is no possibility of using provided abstractions and APIs for module orchestration and execution. Another challenge are limited hardware resources which prevent the deployment of many existing operating system kernels. Basically there are two types of execution models which can be implemented in micro-controllers:

\begin{figure}[H]
    \centering
    \import{images/}{sequential_execution.pdf_tex}
    \caption{Sequential execution model}
\end{figure}

\minisec{Sequential execution}
This type sequentially executes all modules starting from the first module until the last one. Once the last module ends the execution starts again from the first module. It is a very simple model that does not need any operating system support or frameworks. It can be simply implemented as a sequence of function calls inside an infinite loop as shown in algorithm \ref{alg:sequential execution}.

\begin{algorithm}[H]
\caption{Sequential model algorithm}
\label{alg:sequential execution}
\begin{algorithmic}
\WHILE{$true$}
\STATE $module_1$
\STATE $module_2$
\STATE $...$
\STATE $module_n$
\ENDWHILE
\end{algorithmic}
\end{algorithm}

There is one challenge that comes with this type of execution model. That is that only one module can execute at a time due to its sequential nature. If a module i.e. waits for an external resource to provide data it must not block the execution of the main loop until the external resources becomes ready. This would prevent the execution of the other modules. The classic solution to this problem is the introduction of states in modules. Module states can be implemented as classical Finite State Machines (\cite{booth}).

\minisec{}
If we take the example from above about waiting for external resources a finite state machine for modules can be modeled as shown in figure \ref{fig:statemachine}.
\begin{figure}[H]
\centering
\import{images/}{state_machine.pdf_tex}
\caption{State Machine for a module}
\label{fig:statemachine}
\end{figure}

State machine models can be implemented using \texttt{if} or \texttt{case} statements which is shown in algorithm \ref{alg:statemachine}. The nice side effect of a state machine based implementation is the non-blocking nature of the module execution. Take for instance the execution of state 1 "Waiting" as shown in figure \ref{fig:statemachine}. The CPU only needs to execute as many instructions as are necessary to check if the awaited resource is available. If the resource is not available the execution returns to the main loop and the next module (together with its state machine) is being executed.

\begin{algorithm}[H]
\caption{State machine algorithm}
\label{alg:statemachine}
\begin{algorithmic}
\IF{state is WAITING}
    \IF{resource is available}
        \STATE set state to PROCESSING
    \ELSE
        \STATE return
    \ENDIF
\ELSIF{state is PROCESSING}
    \STATE process data
    \STATE set state to WAITING
\ENDIF
\end{algorithmic}
\end{algorithm}

This implementation emulates a concurrent execution of modules. The context switch between module executions is being done by the modules themselves (using self-interruption) and no external scheduler is involved. This form of concurrent behavior can therefore be described as a non-preemptive or cooperative multi-tasking between modules. The predecessor thesis \cite{korniowski} implementation heavily used the described state machine algorithm although the model theory behind the implementation was not being mentioned in the thesis. Listing \ref{lst:korniowski_main} shows the main function of the predecessor thesis implementation.

\begin{lstlisting}[label=lst:korniowski_main,caption=main function implementation in \cite{korniowski}]
382 while(0x01)
383 {
384         if(uartInterrupt == ON) // got a character from RS232
385 +---- 44 lines: 
429 
430 
431         // --- RECEIVE A DATAGRAM ---
432 
433         else if((datagramReceived = datagramReceive(...)) 
                    && netState > 0)     
434 +----182 lines: 
616 
617     
618         else if(helloTime) // prepare periodic Hello message
619 +---- 19 lines: 
638 
639 
640     
641         // --- SEND A DATAGRAM ---
642     
643         if(datagramReady && netState > 0)
644 +----  8 lines: 
652 
653 }
\end{lstlisting}

A couple of problems arise from the existing implementation. First of all listing \ref{lst:korniowski_main} reveals the following modules:

\begin{itemize}
\item UART Module
\item Datagram Receiver Module
\item Hello Message Sender Module
\item Datagram Sender Module
\end{itemize}

Which module is being executed depends on the state of the main module being represented by the main function. The state of the main module on the other hand depends directly from the state of the submodules. The main module therefore acts more like a controller of the submodules and takes away the responsibility of the submodule's state management. Furthermore the main function is very long and complex (271 lines of code). The lack of a clear separation of module responsibility and conformance the state machine theory led to a completely new implementation as show in listing \ref{lst:urbaniak_main}.

\begin{lstlisting}[label=lst:urbaniak_main,caption=main function implementation]
 95 while (true) {
 96   shell();
 97   batman_thread();
 98   rx_thread();
 99   uart_tx_thread();
100   watchdog();
101   timer_thread();
102 }
\end{lstlisting}

The new implementation makes it very clear which modules are being executed sequentially. Furthermore the main function does not act as a controller but rather leaves the state management in the module's responsibility.

\minisec{}
There is a problem though in state machine based implementations and that is the rapidly growing complexity. This problem is called "state explosion problem" and has even a exponential behavior as shown in \cite{katoen}. The equation \ref{eq:state explosion} shows that the number of states is dependent on the number of program locations, the number of used variables and their dimensions.

\begin{equation}
\label{eq:state explosion}
\#states = \abs{\#locations} \cdot \prod_{variable \: x} \: \abs{dom(x)}
\end{equation}

This equation shows that for instance a program having 10 locations and only 3 boolean variables already has 80 different states. Although this equation might not apply exactly to state machine based implementations it underlines the practical experience of big state-machine based implementations. The alternative to state-machine based applications are thread or process based implementations using the concurrent execution model as shown below.

\begin{figure}[H]
\centering
\import{images/}{concurrent_execution.pdf_tex}
\caption[]{Concurrent execution model}
\end{figure}

\minisec{Concurrent execution}
This execution model executes modules concurrently. Instead of heaving an infinite main loop that iterates sequentially over all modules the main function only initializes and launches concurrent modules. Each module runs in isolation and can have its own main loop or terminate immediately. This model requires support from an existing operating system. An existing framework or API provides the necessary abstraction to create new concurrent modules. In terms of operating systems two abstractions are wide-spread for concurrently running software modules:

\begin{itemize}
\item \textbf{Processes}: Processes are usually considered as seperately concurrently running programs. Usually each process owns its own memory context and communication with other processes happens through abstractions like pipes or shared memory.
\item \textbf{Threads}: Threads are concurrently running code parts from the same program. The initial program is considered to run in its own "main thread". Other threads can be started from the main thread. Threads also do run in isolation to each other. Each thread has its own stack. Communication with other threads happens through shared memory provided by static data or the heap.
\end{itemize}

Processes as well as threads are widely known concepts in classical desktop operating systems. In the area of micro-controllers these concepts also are implemented in very lightweight implementations. Many solutions exist like:

\begin{enumerate}
\item FreeRTOS (http://www.freertos.org)
\item TinyOS (http://www.tinyos.net)
\item Atomthreads (http://http://atomthreads.com)
\item Nut/OS (http://www.ethernut.de/en/firmware/nutos.html)
\item BeRTOS (http://www.bertos.org)
\end{enumerate}

The above solutions have chosen different names for threads or processes (some call them "tasks") but essentially they all share the same concept of the concurrent execution model and will be referred to as concurrent modules from now on. Algorithm \ref{alg:threads} shows the pseudo-code that initializes concurrent modules. One can see that in contrast to the sequential execution model the main loop actually does nothing.

\begin{algorithm}[H]
\caption{Concurrent model initialization}
\label{alg:threads}
\begin{algorithmic}

\STATE start $module_1$
\STATE start $module_2$
\STATE ...
\STATE start $module_n$

\WHILE{true}
    \STATE // no operation
\ENDWHILE
\end{algorithmic}
\end{algorithm}

But how does a context switch happen between concurrent modules? Two methodologies exist:

\begin{itemize}
\item \textbf{Cooperative}: The concurrent modules by themselves return the control to a scheduler which then delegates the control to a different module. Which concurrent module gets control is often based on priorities which are controlled by the scheduler.
\item \textbf{Preemptive}: Here the concurrent modules do not have control about how and when they get interrupted. It can happen anytime during the execution. Again the context switch between concurrent modules is often handled using priorities in the scheduler.
\end{itemize}

Nearly all existing solutions have one feature in common. That is that every thread has its own seperate stack memory space. This is necessary in order to be able to run the same block of code (for instance a function) in multiple thread instances. On the other hand threads are being executed in the same memory context so sharing data between threads is possible by using the heap or static memory. All of the above mentioned frameworks provide common abstractions which are needed in thread based implementations:

\begin{itemize}
\item Semaphores
\item Mutexes
\item Yielding
\end{itemize}

\minisec{}
In contrast to state machine based or sequential based concurrency thread based implementation can be expressed in very linear algorithms using the above mentioned abstractions. Take for instance the state-machine based algorithm \ref{alg:statemachine}. This could be translated into a linear thread-based algorithm as shown in \ref{alg:thread}.

\begin{algorithm}[H]
\caption{Thread based algorithm}
\label{alg:thread}
\begin{algorithmic}
\WHILE{true}
    \STATE wait for resource mutex
    \STATE process data
    \STATE release resource mutex
\ENDWHILE
\end{algorithmic}
\end{algorithm}

One can easily see that the thread-based algorithm \ref{alg:thread} is much more expressive than the state-machine based algorithm \ref{alg:statemachine}.

\minisec{}
Together with the necessity of having a scheduler these solutions can be considered as heavy-weight. The scheduler consumes additional CPU cycles and the seperate stack memory space per thread consumes additional memory which is very scarce in embedded micro-controller systems.

\minisec{}
Although ususally a concurrent execution model must be provided in form of an API or an existing kernel there is one exception in the context embedded micro-controllers and that are ISRs (Interrupt Service Routines). Interrupt service routines behave very much like preemptive concurrent modules with highest priority.

\begin{figure}[H]
\centering
\import{images/}{isr.pdf_tex}
\caption[]{Illustration of an Interrupt Service Routine}
\end{figure}

The ISR interrupts the main program at any time when an external resource triggers an event and executes the service routine. The scheduler in this case is the CPU itself. There is one caveat with ISRs. When one ISR is being executed no other ISR can be triggered. Therefore it is being considered best practice not to perform intense and long running operations in ISRs.

Alternatives: \\
- Heavyweight: Real Operating System. Enumerate them and compare ... \\
- Lightweight: Thread implementations. Problem: Each thread has its own stack which consumes a lot of memory. \\
- More Lightweight: Protothreads. Best compromise between classical state machines and real threads. \\

\section{Ring Buffers}
\section{Half-Duplex Radio Access (Petri Net)}

\chapter{Network Stack}
\section{Layer 2a: MAC Layer}
\section{Layer 2b: Logical Link Control}
\section{Layer 3: Batman Routing}
\section{Layer 7: Application}

\chapter{Research}
\section{Simulations}
\subsection{Shell}
\subsection{Routing}
\subsection{Radio Transmission}
\section{Mesh evaluation}
\section{Results}

\chapter{Conclusion}


% Anhang (Bibliographie darf im deutschen nicht in den Anhang!)
\bibliography{bib/BibtexDatabase}
\clearpage
% Abbildungs- und Tabellenverzeichnis
\listoffigures
\listoftables
% Anhang
\appendix
% 'Anhang' ins Inhaltsverzeichnis
%\phantomsection
%\addcontentsline{toc}{part}{Anhang}

% \input{content/Z-Anhang-01-Herleitungen}

\chapter*{Danksagung}


\IfDefined{printindex}{\printindex}
\IfDefined{printnomenclature}{\printnomenclature}

%% Dokument ENDE %%%%%%%%%%%%%%%%%%%%%%%%%%%%%%%%%%%%%%%%%%%%%%%%%%%%%%%%%%
\end{document}

